\input{preamble}

% OK, start here.
%
\begin{document}

\title{La géométrie des champs algébriques}


\maketitle

\phantomsection
\label{section-phantom}

\tableofcontents

\section{Introduction}
\label{section-introduction}

\noindent
Ce chapitre étudie quelques propriétés géométriques des champs algébriques.
Les versions initiales des sections \ref{section-multiplicities} et
\ref{section-dimension-of-algebraic-stacks}
ont été rédigées par Matthew Emerton et Toby Gee et se trouvent
sous leur forme originale dans \cite{Emerton-Gee-dim}.







\section{Anneaux versels}
\label{section-versal}

\noindent
Dans cette section, nous explicitons le lien entre les anneaux de déformation
et les anneaux locaux des champs algébriques de type fini sur une base
localement noethérienne.

\begin{situation}
\label{situation-versal}
Ici, $\mathcal{X}$ est un champ algébrique localement de type fini
sur un schéma localement noethérien $S$.
\end{situation}

\noindent
Voici la définition.

\begin{definition}
\label{definition-versal-ring-at-x}
Dans la Situation \ref{situation-versal}, soit $x_0 : \Spec(k) \to \mathcal{X}$
un morphisme, où $k$ est un corps de type fini sur $S$.
Un {\it anneau versel de $\mathcal{X}$ en $x_0$} est
une $S$-algèbre locale noethérienne complète $A$ de corps résiduel $k$,
pour laquelle il existe un objet formel versel
$(A, \xi_n, f_n)$ comme dans Axiomes d'Artin, Définition
\ref{artin-definition-versal-formal-object},
avec $\xi_1 \cong x_0$ (une $2$-isomorphie).
\end{definition}

\noindent
Nous voulons démontrer que les anneaux versels existent et sont uniques à
facteurs lisses près. Nous utiliserons pour cela les catégories de prédéformations de
Axiomes d'Artin, section \ref{artin-section-predeformation-categories}.
Dans notre situation, ce sont toujours des catégories de déformations.

\begin{lemma}
\label{lemma-deformation-category}
Dans la Situation \ref{situation-versal}, soit $x_0 : \Spec(k) \to \mathcal{X}$
un morphisme, où $k$ est un corps de type fini sur $S$.
Alors $\mathcal{F}_{\mathcal{X}, k, x_0}$
est une catégorie de déformations, et $T\mathcal{F}_{\mathcal{X}, k, x_0}$
ainsi que $\text{Inf}(\mathcal{F}_{\mathcal{X}, k, x_0})$
sont des espaces vectoriels sur $k$ de dimension finie.
\end{lemma}

\begin{proof}
Choisissons un ouvert affine $\Spec(\Lambda) \subset S$ tel que
$\Spec(k) \to S$ se factorise par celui-ci.
D'après Axiomes d'Artin, section \ref{artin-section-predeformation-categories},
nous obtenons une catégorie de prédéformations
$\mathcal{F}_{\mathcal{X}, k, x_0}$
sur la catégorie $\mathcal{C}_\Lambda$.
(Comme il est signalé au passage cité, cette catégorie ne dépend que
du morphisme $\Spec(k) \to S$, et non du choix de
$\Lambda$.) D'après Axiomes d'Artin, Lemmes
\ref{artin-lemma-deformation-category} et
\ref{artin-lemma-algebraic-stack-RS}
$\mathcal{F}_{\mathcal{X}, k, x_0}$ est en fait une catégorie de déformations.
D'après Axiomes d'Artin, Lemme \ref{artin-lemma-finite-dimension},
nous constatons que $T\mathcal{F}_{\mathcal{X}, k, x_0}$
et $\text{Inf}(\mathcal{F}_{\mathcal{X}, k, x_0})$
sont des espaces vectoriels sur $k$ de dimension finie.
\end{proof}

\begin{lemma}
\label{lemma-versal-ring}
Dans la Situation \ref{situation-versal}, soit $x_0 : \Spec(k) \to \mathcal{X}$
un morphisme, où $k$ est un corps de type fini sur $S$.
Il existe alors un anneau versel de $\mathcal{X}$ en $x_0$. Si
$A$, $A'$ en sont deux, alors $A \cong A'[[t_1, \ldots, t_r]]$
ou $A' \cong A[[t_1, \ldots, t_r]]$ comme $S$-algèbres
pour un certain $r$.
\end{lemma}

\begin{proof}
L'existence résulte du Lemme \ref{lemma-deformation-category} et de
Théorie formelle des déformations,
Définition \ref{formal-defos-definition-deformation-category} et
Lemmes \ref{formal-defos-lemma-RS-implies-S1-S2} et
\ref{formal-defos-lemma-versal-object-existence}.
D'après le résultat d'unicité de
Théorie formelle des déformations, Lemme \ref{formal-defos-lemma-minimal-versal},
il existe un anneau versel « minimal » $A$ de $\mathcal{X}$ en $x_0$
tel que tout autre anneau versel de $\mathcal{X}$ en $x_0$ soit
isomorphe à $A[[t_1, \ldots, t_r]]$ pour un certain $r$.
Ceci entraîne manifestement la seconde assertion.
\end{proof}

\begin{lemma}
\label{lemma-versal-ring-field-extension}
Dans la Situation \ref{situation-versal}, soit $x_0 : \Spec(k) \to \mathcal{X}$
un morphisme, où $k$ est un corps de type fini sur $S$.
Soit $l/k$ une extension finie de corps, et notons
$x_{l, 0} : \Spec(l) \to \mathcal{X}$ le morphisme induit.
Étant donné un anneau versel $A$ de $\mathcal{X}$ en $x_0$, il existe
un anneau versel $A'$ de $\mathcal{X}$ en $x_{l, 0}$ tel
qu'il existe un homomorphisme de $S$-algèbres $A \to A'$ induisant
l'extension de corps donnée $l/k$ et formellement lisse pour la topologie $\mathfrak m_{A'}$-adique
considérée.
\end{lemma}

\begin{proof}
Cela résulte immédiatement de
Axiomes d'Artin, Lemme \ref{artin-lemma-change-of-field},
et de
Théorie formelle des déformations, Lemme
\ref{formal-defos-lemma-change-of-field-versal-ring}.
(Nous utilisons aussi le fait que $\mathcal{X}$ satisfait (RS), d'après
Axiomes d'Artin, Lemme \ref{artin-lemma-algebraic-stack-RS}.)
\end{proof}

\begin{lemma}
\label{lemma-compare-versal-ring-completion}
Dans la Situation \ref{situation-versal}, soit $x : U \to \mathcal{X}$ un
morphisme, où $U$ est un schéma localement de type fini sur $S$.
Soit $u_0 \in U$ un point de type fini.
Posons $k = \kappa(u_0)$ et notons $x_0 : \Spec(k) \to \mathcal{X}$
l'application induite. Les conditions suivantes sont équivalentes :
\begin{enumerate}
\item $x$ est versel en $u_0$
(Axiomes d'Artin, Définition \ref{artin-definition-versal}),
\item $\hat x : \mathcal{F}_{U, k, u_0} \to \mathcal{F}_{\mathcal{X}, k, x_0}$
est lisse,
\item l'objet formel associé à
$x|_{\Spec(\mathcal{O}_{U, u_0}^\wedge)}$ est versel, et
\item il existe un voisinage ouvert $U' \subset U$ de $x$ tel que
$x|_{U'} : U' \to \mathcal{X}$ soit lisse.
\end{enumerate}
De plus, dans ce cas, le complété $\mathcal{O}_{U, u_0}^\wedge$
est un anneau versel de $\mathcal{X}$ en $x_0$.
\end{lemma}

\begin{proof}
Puisque $U \to S$ est localement de type fini (comme composé de tels morphismes),
nous voyons que $\Spec(k) \to S$ est de type fini (encore comme composé).
L'énoncé a donc un sens. L'équivalence de (1) et (2)
est la définition du fait que $x$ est versel en $u_0$.
L'équivalence de (1) et (3) est
Axiomes d'Artin, Lemme \ref{artin-lemma-versality-matches}.
Ainsi, (1), (2) et (3) sont équivalentes.

\medskip\noindent
Si $x|_{U'}$ est lisse, alors le foncteur
$\hat x : \mathcal{F}_{U, k, u_0} \to \mathcal{F}_{\mathcal{X}, k, x_0}$
est lisse d'après Axiomes d'Artin, Lemme
\ref{artin-lemma-formally-smooth-on-deformation-categories}.
Ainsi, (4) entraîne (1), (2) et (3).
Réciproquement, supposons que $x$ soit versel en $u_0$.
Choisissons un morphisme lisse surjectif $y : V \to \mathcal{X}$, où $V$
est un schéma. Posons $Z = V \times_\mathcal{X} U$ et choisissons un
point de type fini $z_0 \in |Z|$ au-dessus de $u_0$ (cela est possible d'après
Morphismes d'espaces, Lemme
\ref{spaces-morphisms-lemma-finite-type-points-surjective-morphism}).
D'après Axiomes d'Artin, Lemme \ref{artin-lemma-base-change-versal},
le morphisme $Z \to V$ est lisse en $z_0$.
Par définition, nous pouvons trouver un voisinage ouvert $W \subset Z$
de $z_0$ tel que $W \to V$ soit lisse. Puisque $Z \to U$ est ouvert,
soit $U' \subset U$ l'image de $W$. Nous voyons alors que
$U' \to \mathcal{X}$ est lisse par notre définition des morphismes lisses
de champs.

\medskip\noindent
La dernière assertion résulte des définitions, puisque
$\mathcal{O}_{U, u_0}^\wedge$
proreprésente $\mathcal{F}_{U, k, u_0}$.
\end{proof}

\begin{lemma}
\label{lemma-characterize-smoothness}
Dans la Situation \ref{situation-versal}. Soit $x_0 : \Spec(k) \to \mathcal{X}$
un morphisme tel que $\Spec(k) \to S$ soit de type fini, d'image $s$.
Soit $A$ un anneau versel de $\mathcal{X}$ en $x_0$. Les conditions suivantes
sont équivalentes :
\begin{enumerate}
\item $x_0$ appartient au lieu lisse de $\mathcal{X} \to S$
(Morphismes de champs, Lemme \ref{stacks-morphisms-lemma-where-smooth}),
\item $\mathcal{O}_{S, s} \to A$ est formellement lisse pour la
topologie $\mathfrak m_A$-adique, et
\item $\mathcal{F}_{\mathcal{X}, k, x_0}$ est sans obstruction.
\end{enumerate}
\end{lemma}

\begin{proof}
L'équivalence de (2) et (3) résulte immédiatement de
Théorie formelle des déformations, Lemme
\ref{formal-defos-lemma-smooth-power-series-classical}.

\medskip\noindent
Notons que $\mathcal{O}_{S, s} \to A$ est formellement lisse pour la
topologie $\mathfrak m_A$-adique si et seulement si
$\mathcal{O}_{S, s} \to A' = A[[t_1, \ldots, t_r]]$
est formellement lisse pour la topologie $\mathfrak m_{A'}$-adique.
Par conséquent, (2) ne dépend pas du choix
de l'anneau versel, d'après le Lemme \ref{lemma-versal-ring}.
Ensuite, soit $l/k$ une extension finie et choisissons
$A \to A'$ comme dans le Lemme \ref{lemma-versal-ring-field-extension}.
Si $\mathcal{O}_{S, s} \to A$ est formellement lisse pour la
topologie $\mathfrak m_A$-adique, alors $\mathcal{O}_{S, s} \to A'$
est formellement lisse pour la topologie $\mathfrak m_{A'}$-adique, voir
Compléments d'algèbre, Lemme \ref{more-algebra-lemma-compose-formally-smooth}.
Réciproquement, si $\mathcal{O}_{S, s} \to A'$
est formellement lisse pour la topologie $\mathfrak m_{A'}$-adique,
alors $\mathcal{O}_{S, s}^\wedge \to A'$ et $A \to A'$ sont réguliers
(Compléments d'algèbre, Proposition \ref{more-algebra-proposition-fs-regular}),
donc $\mathcal{O}_{S, s}^\wedge \to A$ est régulier
(Compléments d'algèbre, Lemme \ref{more-algebra-lemma-regular-permanence}),
et donc $\mathcal{O}_{S, s} \to A$ est formellement lisse pour la
topologie $\mathfrak m_A$-adique (même lemme que précédemment).
Ainsi, l'équivalence de (2) et (1) vaut pour
$k$ et $x_0$ si et seulement si elle vaut pour $l$ et $x_{0, l}$.

\medskip\noindent
Choisissons un schéma $U$ et un morphisme lisse $U \to \mathcal{X}$ tels
que $\Spec(k) \times_\mathcal{X} U$ soit non vide. Choisissons une extension
finie $l/k$ et un point $w_0 : \Spec(l) \to \Spec(k) \times_\mathcal{X} U$.
Soit $u_0 \in U$ l'image de $w_0$.
Nous pouvons appliquer ce qui précède à $l/k$ et à $l/\kappa(u_0)$
pour voir que nous pouvons nous ramener à $u_0$. Nous pouvons donc supposer
$A = \mathcal{O}_{U, u_0}^\wedge$, voir le
Lemme \ref{lemma-compare-versal-ring-completion}.
Observons que $x_0$ appartient au lieu lisse de $\mathcal{X} \to S$
si et seulement si $u_0$ appartient au lieu lisse de $U \to S$, voir par exemple
Morphismes de champs, Lemme \ref{stacks-morphisms-lemma-where-smooth}.
L'équivalence de (1) et (2) résulte donc de
Compléments d'algèbre, Lemme \ref{more-algebra-lemma-formally-smooth-finite-type}.
\end{proof}

\noindent
Rappelons une conséquence de l'approximation d'Artin.

\begin{lemma}
\label{lemma-Artin-approximation-by-smooth-morphism}
Dans la Situation \ref{situation-versal}. Soit $x_0 : \Spec(k) \to \mathcal{X}$
un morphisme tel que $\Spec(k) \to S$ soit de type fini, d'image $s$.
Soit $A$ un anneau versel de $\mathcal{X}$ en $x_0$.
Si $\mathcal{O}_{S, s}$ est un anneau G, nous pouvons trouver un morphisme lisse
$U \to \mathcal{X}$ dont la source est un schéma, et un point
$u_0 \in U$ de corps résiduel $k$, tels que
\begin{enumerate}
\item $\Spec(k) \to U \to \mathcal{X}$ coïncide avec le morphisme donné $x_0$,
\item il existe un isomorphisme $\mathcal{O}_{U, u_0}^\wedge \cong A$.
\end{enumerate}
\end{lemma}

\begin{proof}
Soit $(\xi_n, f_n)$ l'objet formel versel sur $A$.
D'après Axiomes d'Artin, Lemme \ref{artin-lemma-effective},
nous savons que $\xi = (A, \xi_n, f_n)$ est effectif.
Par hypothèse, $\mathcal{X}$ est localement de présentation finie sur $S$
(utiliser Morphismes de champs, Lemme
\ref{stacks-morphisms-lemma-noetherian-finite-type-finite-presentation}),
et préserve donc les limites d'après Limites de champs, Proposition
\ref{stacks-limits-proposition-characterize-locally-finite-presentation}.
Ainsi, l'approximation d'Artin, sous la forme de
Axiomes d'Artin, Lemme \ref{artin-lemma-approximate-versal},
montre que nous pouvons trouver un morphisme $U \to \mathcal{X}$ dont
la source est un $S$-schéma de type fini, contenant un point $u_0 \in U$
de corps résiduel $k$ satisfaisant (1) et (2), tel que $U \to \mathcal{X}$
soit versel en $u_0$. D'après le Lemme \ref{lemma-compare-versal-ring-completion},
quitte à rétrécir $U$, nous pouvons supposer que $U \to \mathcal{X}$ est lisse.
\end{proof}

\begin{remark}[Amélioration des anneaux versels]
\label{remark-upgrade}
Dans la Situation \ref{situation-versal}, soit $x_0 : \Spec(k) \to \mathcal{X}$
un morphisme, où $k$ est un corps de type fini sur $S$.
Soit $A$ un anneau versel de $\mathcal{X}$ en $x_0$. D'après Axiomes d'Artin,
Lemme \ref{artin-lemma-effective}, notre objet formel versel
provient en fait d'un morphisme
$$
\Spec(A) \longrightarrow \mathcal{X}
$$
sur $S$. De plus, chacun des résultats précédents peut être renforcé de façon
à être compatible avec ce morphisme. En voici la liste :
\begin{enumerate}
\item dans le Lemme \ref{lemma-versal-ring}, l'isomorphisme
$A \cong A'[[t_1, \ldots, t_r]]$ ou $A' \cong A[[t_1, \ldots, t_r]]$
peut être choisi compatible avec ces morphismes,
\item dans le Lemme \ref{lemma-versal-ring-field-extension},
l'homomorphisme $A \to A'$ peut être choisi compatible avec ces morphismes,
\item dans le Lemme \ref{lemma-compare-versal-ring-completion},
le morphisme $\Spec(\mathcal{O}_{U, u_0}^\wedge) \to \mathcal{X}$
est le composé de l'application canonique
$\Spec(\mathcal{O}_{U, u_0}^\wedge) \to U$ et de l'application donnée
$U \to \mathcal{X}$,
\item dans le Lemme \ref{lemma-Artin-approximation-by-smooth-morphism},
l'isomorphisme $\mathcal{O}_{U, u_0}^\wedge \cong A$ peut
être choisi de telle sorte que $\Spec(A) \to \mathcal{X}$ corresponde à l'application canonique
de l'alinéa précédent.
\end{enumerate}
Dans chaque cas, l'assertion résulte du fait que nos applications sont
compatibles avec les éléments formels versels; notons toutefois que les
diagrammes sous-entendus ne sont $2$-commutatifs qu'au choix (non canonique)
d'une $2$-flèche près. Cela signifie néanmoins que l'application $A' \to A$
ou $A \to A'$ sous-entendue dans (1) est bien définie à homotopie formelle près, voir
Théorie formelle des déformations, Lemme
\ref{formal-defos-lemma-versal-unique-up-to-homotopy}.
\end{remark}

\begin{lemma}
\label{lemma-versal-ring-flat}
Dans la Situation \ref{situation-versal}, soit $x_0 : \Spec(k) \to \mathcal{X}$
un morphisme, où $k$ est un corps de type fini sur $S$.
Soit $A$ un anneau versel de $\mathcal{X}$ en $x_0$.
Alors le morphisme $\Spec(A) \to \mathcal{X}$ de la
Remarque \ref{remark-upgrade} est plat.
\end{lemma}

\begin{proof}
Si l'anneau local de $S$ au point image est un anneau G, cela
résulte immédiatement du
Lemme \ref{lemma-Artin-approximation-by-smooth-morphism}
et du fait que l'application d'un anneau local noethérien dans son
complété est plate. Dans le cas général, nous procédons comme suit.

\medskip\noindent
Étape I. Si $A$ et $A'$ sont deux anneaux versels de $\mathcal{X}$ en $x_0$,
le résultat est vrai pour $A$ si et seulement s'il l'est pour $A'$.
En effet, quitte à échanger $A$ et $A'$, nous pouvons supposer qu'il existe
une application formellement lisse $\varphi : A \to A'$ telle que le composé
$$
\Spec(A') \to \Spec(A) \to \mathcal{X}
$$
soit le morphisme $\Spec(A') \to \mathcal{X}$, voir le
Lemme \ref{lemma-versal-ring} et la
Remarque \ref{remark-upgrade}.
Puisque $A \to A'$ est fidèlement plat, l'équivalence résulte de
Morphismes de champs, Lemmes \ref{stacks-morphisms-lemma-composition-flat} et
\ref{stacks-morphisms-lemma-flat-permanence}.

\medskip\noindent
Étape II. Soit $l/k$ une extension finie de corps. Soit
$x_{l, 0} : \Spec(l) \to \mathcal{X}$ le morphisme induit.
Soit $A$ un anneau versel de $\mathcal{X}$ en $x_0$ et
soit $A \to A'$ comme dans le Lemme \ref{lemma-versal-ring-field-extension}.
À nouveau, le composé
$$
\Spec(A') \to \Spec(A) \to \mathcal{X}
$$
est le morphisme $\Spec(A') \to \mathcal{X}$, voir la Remarque \ref{remark-upgrade}.
En raisonnant comme précédemment et en utilisant l'étape I pour voir que le choix des anneaux versels
est sans incidence, nous voyons que le lemme vaut pour $x_0$ si et seulement
s'il vaut pour $x_{l, 0}$.

\medskip\noindent
Étape III. Choisissons un schéma $U$ et un morphisme lisse surjectif
$U \to \mathcal{X}$. Nous pouvons alors choisir un point de type fini
$z_0$ de $Z = U \times_\mathcal{X} x_0$ (c'est un espace algébrique
non vide). Soit $u_0 \in U$ l'image de $z_0$ dans $U$.
Choisissons un schéma et une application étale surjective $W \to Z$
tels que $z_0$ soit l'image d'un point fermé $w_0 \in W$
(voir Morphismes d'espaces, section
\ref{spaces-morphisms-section-points-finite-type}).
Puisque $W \to \Spec(k)$ et $W \to U$ sont de type fini,
nous voyons que $\kappa(w_0)/k$ et $\kappa(w_0)/\kappa(u_0)$
sont des extensions finies de corps
(voir Morphismes, section \ref{morphisms-section-points-finite-type}).
En appliquant deux fois l'étape II, nous pouvons remplacer $x_0$ par
$u_0 \to U \to \mathcal{X}$. Nous voyons alors que notre morphisme est le composé
$$
\Spec(\mathcal{O}_{U, u_0}^\wedge) \to U \to \mathcal{X}
$$
La première flèche est plate parce que les complétions d'anneaux locaux noethériens
sont plates (Algèbre, Lemme \ref{algebra-lemma-completion-flat}),
et la seconde flèche est plate puisque
les morphismes lisses sont plats. Le composé est plat car la composition
préserve la platitude.
\end{proof}

\begin{remark}
\label{remark-groupoid-defo}
Dans la Situation \ref{situation-versal}, soit $x_0 : \Spec(k) \to \mathcal{X}$
un morphisme, où $k$ est un corps de type fini sur $S$.
D'après le Lemme \ref{lemma-deformation-category} et
Théorie formelle des déformations, Théorème
\ref{formal-defos-theorem-presentation-deformation-groupoid}
nous savons que $\mathcal{F}_{\mathcal{X}, k, x_0}$ admet une
présentation par un groupoïde lisse proreprésentable en
foncteurs sur $\mathcal{C}_\Lambda$.
En explicitant les définitions, cela signifie que nous pouvons choisir
\begin{enumerate}
\item une $\Lambda$-algèbre locale noethérienne complète $A$
de corps résiduel $k$, et un objet formel versel $\xi$
de $\mathcal{F}_{\mathcal{X}, k, x_0}$ sur $A$,
\item une $\Lambda$-algèbre locale noethérienne complète $B$
de corps résiduel $k$, et un isomorphisme
$$
\underline{B}|_{\mathcal{C}_\Lambda}
\longrightarrow
\underline{A}|_{\mathcal{C}_\Lambda}
\times_{\underline{\xi}, \mathcal{F}_{\mathcal{X}, k, x_0}, \underline{\xi}}
\underline{A}|_{\mathcal{C}_\Lambda}
$$
\end{enumerate}
Les projections correspondent à des applications formellement lisses
$t : A \to B$ et $s : A \to B$ (car $\xi$ est versel).
Il existe une application $c : B \to B \widehat{\otimes}_{s, A, t} B$
qui fait de $(A, B, s, t, c)$ un cogroupoïde dans la catégorie
des $\Lambda$-algèbres locales noethériennes complètes de corps résiduel $k$
(sur les foncteurs proreprésentables, cette application est construite dans
Théorie formelle des déformations, Lemme
\ref{formal-defos-lemma-presentation-construction}).
Enfin, le théorème cité nous dit que $\xi$ induit
une équivalence
$$
[\underline{A}|_{\mathcal{C}_\Lambda} / \underline{B}|_{\mathcal{C}_\Lambda}]
\longrightarrow
\mathcal{F}_{\mathcal{X}, k, x_0}
$$
de groupoïdes cofibrés sur $\mathcal{C}_\Lambda$. En fait, nous obtenons aussi
une équivalence
$$
[\underline{A}/\underline{B}]
\longrightarrow
\widehat{\mathcal{F}}_{\mathcal{X}, k, x_0}
$$
de groupoïdes cofibrés sur la catégorie complétée
$\widehat{\mathcal{C}}_\Lambda$ (voir la discussion dans
Théorie formelle des déformations, section
\ref{formal-defos-section-prorepresentable-groupoids-in-functors}
qui explique pourquoi cela fonctionne). Bien entendu, $A$ est un anneau versel de
$\mathcal{X}$ en $x_0$.
\end{remark}








\section{Multiplicités des composantes des champs algébriques}
\label{section-multiplicities}

\noindent
Si $X$ est un schéma localement noethérien, nous pouvons écrire $X$ (considéré
simplement comme espace topologique) comme réunion de composantes irréductibles,
disons $X = \bigcup T_i$. Chaque composante irréductible est l'adhérence d'un
unique point générique $\xi_i$, et l'anneau local $\mathcal O_{X,\xi_i}$
est un anneau artinien local. Nous pouvons définir la {\it multiplicité de $X$ le long de $T_i$}
ou la {\it multiplicité de $T_i$ dans $X$} par
$$
m_{T_i, X} = \text{length}_{\mathcal O_{X, \xi_i}} \mathcal O_{X, \xi_i}
$$
Autrement dit, c'est la longueur de l'anneau artinien local. Comparer
avec
Homologie de Chow, section \ref{chow-section-cycle-of-closed-subscheme}.

\medskip\noindent
Notre but est ici de généraliser cette définition aux champs algébriques
localement noethériens. Si $\mathcal{X}$ est un champ,
alors son espace topologique $|\mathcal{X}|$
(voir Propriétés des champs, Définition
\ref{stacks-properties-definition-topological-space})
est localement noethérien
(Morphismes de champs, Lemme \ref{stacks-morphisms-lemma-Noetherian-topology}).
Les composantes irréductibles de $|\mathcal{X}|$ sont parfois
appelées les composantes irréductibles de $\mathcal{X}$.
Si $\mathcal{X}$ est quasi-séparé, alors $|\mathcal{X}|$
est sobre (Morphismes de champs, Lemme
\ref{stacks-morphisms-lemma-sober-qs}),
mais il ne l'est pas nécessairement dans le cas
non quasi-séparé. Considérons par exemple l'espace algébrique non quasi-séparé
$X = \mathbf{A}^1_\mathbf{C}/\mathbf{Z}$.
En outre, il n'existe pas de faisceau structural
sur $|\mathcal{X}|$ dont les fibres puissent servir à définir les multiplicités.

\begin{lemma}
\label{lemma-map-of-components}
Soit $f : U \to \mathcal{X}$ un morphisme lisse d'un schéma
vers un champ algébrique localement noethérien. L'adhérence de l'image de toute
composante irréductible de $|U|$ est une composante irréductible de $|\mathcal{X}|$.
Si $U \to \mathcal{X}$ est surjectif, toutes les composantes irréductibles de
$|\mathcal{X}|$ s'obtiennent ainsi.
\end{lemma}

\begin{proof}
L'application $|U| \to |\mathcal{X}|$ est continue et ouverte d'après
Propriétés des champs, Lemme \ref{stacks-properties-lemma-topology-points}.
Soit $T \subset |U|$ une composante irréductible. Puisque $U$ est localement
noethérien, nous pouvons trouver un ouvert affine non vide $W \subset U$ contenu dans $T$.
Alors $f(T) \subset |\mathcal{X}|$ est irréductible et contient le
sous-ensemble ouvert non vide $f(W)$. Ainsi, l'adhérence de $f(T)$ est irréductible et
contient un ouvert non vide. Il s'ensuit que cette adhérence est une composante
irréductible.

\medskip\noindent
Supposons $U \to \mathcal{X}$ surjectif et soit $Z \subset |\mathcal{X}|$
une composante irréductible. Choisissons un sous-ensemble ouvert noethérien $V$
de $|\mathcal{X}|$ rencontrant $Z$. Après avoir retiré les autres composantes
irréductibles de $V$, nous pouvons supposer que $V \subset Z$.
Prenons une composante irréductible de l'ouvert non vide
$f^{-1}(V) \subset |U|$ et soit $T \subset |U|$ son adhérence.
C'est une composante irréductible de $|U|$, et l'adhérence de $f(T)$
doit coïncider avec $Z$ par notre choix de $T$.
\end{proof}

\noindent
Le lemme précédent s'applique en particulier aux morphismes lisses
entre schémas localement noethériens. Ce cas particulier est
invoqué implicitement dans l'énoncé du lemme suivant.

\begin{lemma}
\label{lemma-multiplicities}
Soit $U \to X$ un morphisme lisse de schémas localement noethériens.
Soit $T'$ une composante irréductible de $U$. Soit $T$ la
composante irréductible de $X$ obtenue comme adhérence de
l'image de $T'$. Alors $m_{T', U} = m_{T, X}$.
\end{lemma}

\begin{proof}
Notons $\xi'$ le point générique de $T'$, et $\xi$ le
point générique de $T$. Soient $A = \mathcal{O}_{X, \xi}$ et
$B = \mathcal{O}_{U, \xi'}$. Nous devons montrer que
$\text{length}_A A = \text{length}_B B$. Puisque
$A \to B$ est un homomorphisme local plat d'anneaux
(car les morphismes lisses sont plats), nous avons
$$
\text{length}_A(A) \text{length}_B(B/\mathfrak m_A B) =
\text{length}_B(B)
$$
d'après Algèbre, Lemme \ref{algebra-lemma-pullback-module}. Il suffit donc
de montrer que $\mathfrak m_A B = \mathfrak m_B$, ou, de manière équivalente, que
$B/\mathfrak m_A B$ est réduit. Puisque $U \to X$ est lisse,
son changement de base $U_{\xi} \to \Spec \kappa(\xi)$ l'est aussi. Comme $U_{\xi}$ est un
schéma lisse sur un corps, il est réduit, et son anneau local
en tout point l'est donc aussi
(Variétés, Lemme \ref{varieties-lemma-smooth-geometrically-normal}).
En particulier,
$$
B/\mathfrak m_A B =
\mathcal{O}_{U, \xi'}/\mathfrak m_{X, \xi}\mathcal{O}_{U, \xi'} =
\mathcal{O}_{U_\xi, \xi'}
$$
est réduit, comme voulu.
\end{proof}

\noindent
À l'aide de ce résultat, nous pouvons montrer qu'il existe une bonne notion
de multiplicité en nous plaçant localement pour la topologie lisse.

\begin{lemma}
\label{lemma-multiplicity}
Soient $U_1 \to \mathcal{X}$ et $U_2 \to \mathcal{X}$ deux morphismes lisses
de schémas vers un champ algébrique localement noethérien $\mathcal{X}$.
Soient $T_1'$ et $T_2'$ des composantes irréductibles de $|U_1|$
et $|U_2|$, respectivement. Supposons que les adhérences des images de
$T_1'$ et $T_2'$ soient la même composante irréductible $T$ de $|\mathcal{X}|$.
Alors $m_{T_1', U_1} = m_{T_2', U_2}$.
\end{lemma}

\begin{proof}
Soient $V_1$ et $V_2$ des sous-ensembles denses de $T_1'$ et $T'_2$, respectivement,
qui soient ouverts dans $U_1$ et $U_2$, respectivement (voir la preuve du
Lemme \ref{lemma-map-of-components}).
Les images de $|V_1|$ et $|V_2|$ dans $|\mathcal{X}|$ sont des sous-ensembles ouverts
non vides de l'ensemble irréductible $T$, et ont donc une intersection
non vide. D'après
Propriétés des champs, Lemme \ref{stacks-properties-lemma-points-cartesian},
l'application $|V_1 \times_\mathcal{X} V_2| \to |V_1| \times_{|\mathcal{X}|} |V_2|$
est surjective. Par conséquent, $V_1 \times_\mathcal{X} V_2$
est un espace algébrique non vide; nous pouvons donc choisir une
surjection étale $V \to V_1 \times_\mathcal{X} V_2$
dont la source est un schéma (non vide).
Si $T'$ est une composante irréductible quelconque de $V$,
alors le Lemme \ref{lemma-map-of-components} montre que l'adhérence de
l'image de $T'$ dans $U_1$ (respectivement dans $U_2$) est égale à $T'_1$
(respectivement à $T'_2$).

\medskip\noindent
En appliquant deux fois le Lemme \ref{lemma-multiplicities}, nous trouvons
que
$$
m_{T_1', U_1} = m_{T', V} = m_{T_2', U_2},
$$
comme voulu.
\end{proof}

\noindent
Nous en avons maintenant fait assez pour montrer que la définition
suivante a un sens.

\begin{definition}
\label{definition-multiplicity}
Soit $\mathcal{X}$ un champ algébrique localement noethérien. Soit
$T \subset |\mathcal{X}|$ une composante irréductible.
La {\it multiplicité} de $T$ dans $\mathcal{X}$ est définie par
$m_{T, \mathcal{X}} = m_{T', U}$, où $f : U \to \mathcal{X}$
est un morphisme lisse d'un schéma et $T' \subset |U|$
est une composante irréductible telle que $f(T') \subset T$.
\end{definition}

\noindent
Ceci est indépendant du choix de $f : U \to \mathcal{X}$
et de celui de la composante irréductible $T'$ envoyée dans $T$,
d'après les Lemmes \ref{lemma-map-of-components} et \ref{lemma-multiplicity}.

\medskip\noindent
Pour terminer, notons qu'il est parfois commode de considérer
une composante irréductible de $\mathcal{X}$ comme un sous-champ fermé.
À cette fin, si $T \subset |\mathcal{X}|$ est une composante irréductible,
nous pouvons considérer l'unique sous-champ fermé réduit
$\mathcal{T} \subset \mathcal{X}$ tel que $|\mathcal{T}| = T$, voir
Propriétés des champs, Définition
\ref{stacks-properties-definition-reduced-induced-stack}.
Si $\mathcal{X}$ est quasi-séparé, alors
une composante irréductible est un champ intègre; voir
Morphismes de champs, section \ref{stacks-morphisms-section-integral-stacks},
pour plus de détails.



\section{Branches formelles et multiplicités}
\label{section-formal-branches}

\noindent
Il sera commode de comparer la notion de multiplicité
d'une composante irréductible donnée par la Définition \ref{definition-multiplicity}
à la notion connexe de multiplicité des composantes irréductibles
(des spectres) des anneaux versels de $\mathcal{X}$ aux points de type fini.

\medskip\noindent
Dans la Situation \ref{situation-versal}, soit $x_0 : \Spec(k) \to \mathcal{X}$
un morphisme, où $k$ est un corps de type fini sur $S$.
Soient $A$, $A'$ des anneaux versels de $\mathcal{X}$ en $x_0$.
Quitte à échanger $A$ et $A'$, nous savons qu'il existe une application
formellement lisse\footnote{Au sens où
$A'$ devient isomorphe à un anneau de séries formelles sur $A$.}
$\varphi : A \to A'$ compatible avec les objets formels versels, voir le
Lemme \ref{lemma-versal-ring} et la
Remarque \ref{remark-upgrade}.
De plus, $\varphi$ est bien définie à homotopie formelle près, voir
Théorie formelle des déformations, Lemme
\ref{formal-defos-lemma-versal-unique-up-to-homotopy}.
En particulier, nous constatons que $\varphi(\mathfrak p)A'$ est un idéal
bien défini de $A'$, d'après
Théorie formelle des déformations, Lemme
\ref{formal-defos-lemma-homotopic-minimal-prime}.
Puisque $A \to A'$ est formellement lisse,
$\varphi(\mathfrak p)A'$ est en fait un idéal premier minimal de $A'$,
et tout idéal premier minimal de $A'$ est de cette forme pour
un unique idéal premier minimal $\mathfrak p \subset A$ (tout cela
se démontre facilement en écrivant $A'$ comme anneau de séries formelles sur $A$).
Ainsi, puisque les idéaux premiers minimaux correspondent
aux composantes irréductibles, la définition suivante a un sens.

\begin{definition}
\label{definition-formal-branches}
Soit $\mathcal{X}$ un champ algébrique localement de type fini
sur un schéma localement noethérien $S$. Soit $x_0 : \Spec(k) \to \mathcal{X}$
un morphisme, où $k$ est un corps de type fini sur $S$.
L'ensemble des {\it branches formelles de $\mathcal{X}$ passant par $x_0$}
est l'ensemble des composantes irréductibles de $\Spec(A)$
pour un choix quelconque d'anneau versel de $\mathcal{X}$ en $x_0$,
ces ensembles étant identifiés pour les différents choix de $A$ par le procédé décrit ci-dessus.
\end{definition}

\noindent
Supposons, dans la situation de la Définition \ref{definition-formal-branches},
qu'une extension finie $l/k$ soit donnée. Posons
$x_{l, 0} : \Spec(l) \to \mathcal{X}$ égal au composé de
$\Spec(l) \to \Spec(k)$ avec $x_0$. Soit
$A \to A'$ comme dans le Lemme \ref{lemma-versal-ring-field-extension}.
Puisque $A \to A'$ est fidèlement plat, le morphisme
$$
\Spec(A') \to \Spec(A)
$$
envoie les (points génériques des) composantes irréductibles sur les
(points génériques des) composantes irréductibles.
Cette application est surjective, mais
elle n'est en général pas bijective.
Autrement dit, nous obtenons une application surjective
$$
\text{branches formelles de }\mathcal{X}\text{ passant par }x_{l, 0}
\longrightarrow
\text{branches formelles de }\mathcal{X}\text{ passant par }x_0
$$
Il se trouve que, si $l/k$ est purement inséparable, alors
l'application est également injective (nous ajouterons ici un énoncé précis
et une preuve si nous en avons un jour besoin).

\begin{lemma}
\label{lemma-branches}
Dans la situation de la Définition \ref{definition-formal-branches},
il existe une surjection canonique de l'ensemble des branches formelles de
$\mathcal{X}$ passant par $x_0$ sur l'ensemble des composantes irréductibles de
$|\mathcal{X}|$ qui contiennent $x_0$ dans $|\mathcal{X}|$.
\end{lemma}

\begin{proof}
Soit $A$ comme dans la Définition \ref{definition-formal-branches}, et
soit $\Spec(A) \to \mathcal{X}$ comme dans la Remarque \ref{remark-upgrade}.
Nous affirmons que le point générique d'une composante irréductible de
$\Spec(A)$ est envoyé sur un point générique d'une composante irréductible
de $|\mathcal{X}|$. Choisissons un schéma $U$ et un morphisme
lisse surjectif $U \to \mathcal{X}$. Considérons le diagramme
$$
\xymatrix{
\Spec(A) \times_\mathcal{X} U \ar[d]_p \ar[r]_-q & U \ar[d]^f \\
\Spec(A) \ar[r]^j & \mathcal{X}
}
$$
D'après le Lemme \ref{lemma-versal-ring-flat}, $j$ est plat.
Par conséquent, $q$ est plat. D'autre part, $f$ est lisse surjectif,
donc $p$ est lisse surjectif. Cela entraîne que tout point générique
$\eta \in \Spec(A)$ d'une composante irréductible est l'image d'un
point de codimension $0$, $\eta'$, de l'espace algébrique
$\Spec(A) \times_\mathcal{X} U$ (voir
Propriétés des espaces, section \ref{spaces-properties-section-generic-points},
pour la notation et l'emploi du théorème de descente dans les anneaux locaux étales).
Puisque $q$ est plat, $q(\eta')$ est un point de codimension $0$ de $U$
(même raisonnement).
Puisque $U$ est un schéma, $q(\eta')$ est le point générique
d'une composante irréductible de $U$. Ainsi, l'adhérence de l'image
de $q(\eta')$ dans $|\mathcal{X}|$ est une composante irréductible
d'après le Lemme \ref{lemma-map-of-components}, comme affirmé.

\medskip\noindent
L'affirmation fournit manifestement un procédé pour définir l'application voulue.
Pour voir qu'elle est surjective, choisissons $u_0 \in U$ s'envoyant sur
$x_0$ dans $|\mathcal{X}|$. Choisissons un voisinage ouvert affine $U' \subset U$
de $u_0$. Quitte à rétrécir $U'$, nous pouvons supposer que toute
composante irréductible de $U'$ passe par $u_0$. Nous pouvons alors
remplacer $\mathcal{X}$ par le sous-champ ouvert correspondant à
l'image de $|U'| \to |\mathcal{X}|$. Nous pouvons donc supposer que $U$ est affine,
possède un point $u_0$ s'envoyant sur $x_0 \in |\mathcal{X}|$,
et que toute composante irréductible de $U$ passe par $u_0$.
D'après Propriétés des champs, Lemme
\ref{stacks-properties-lemma-points-cartesian}
il existe un point $t \in |\Spec(A) \times_\mathcal{X} U|$
s'envoyant sur le point fermé de $\Spec(A)$ et sur $u_0$.
En appliquant le théorème de descente aux homomorphismes locaux plats d'anneaux
$$
A \longrightarrow
\mathcal{O}_{\Spec(A) \times_\mathcal{X} U, \overline{t}}
\longleftarrow \mathcal{O}_{U, u_0}
$$
nous voyons que tout idéal premier minimal de $\mathcal{O}_{U, u_0}$
est l'image d'un idéal premier minimal de l'anneau local du milieu,
et qu'un tel idéal premier minimal s'envoie sur un idéal premier minimal de $A$.
Cela démontre la surjectivité. Certains détails sont omis.
\end{proof}

\noindent
Soit $A$ un anneau local noethérien complet. Les composantes
irréductibles de $\Spec(A)$ ont alors des multiplicités, voir
l'introduction de la section \ref{section-multiplicities}.
Si $A' = A[[t_1, \ldots, t_r]]$, alors le morphisme
$\Spec(A') \to \Spec(A)$ induit une bijection entre les composantes
irréductibles qui préserve les multiplicités (nous omettons la preuve facile).
Ceci, joint à la discussion précédant la
Définition \ref{definition-formal-branches},
montre que la définition suivante a un sens.

\begin{definition}
\label{definition-multiplicity-formal-branches}
Soit $\mathcal{X}$ un champ algébrique localement de type fini
sur un schéma localement noethérien $S$. Soit $x_0 : \Spec(k) \to \mathcal{X}$
un morphisme, où $k$ est un corps de type fini sur $S$.
La {\it multiplicité d'une branche formelle de $\mathcal{X}$ passant par $x_0$}
est la multiplicité de la composante irréductible correspondante de
$\Spec(A)$ pour un choix quelconque d'anneau versel de $\mathcal{X}$ en $x_0$
(voir la discussion ci-dessus).
\end{definition}

\begin{lemma}
\label{lemma-branches-multiplicity}
Soit $\mathcal{X}$ un champ algébrique localement de type fini
sur un schéma localement noethérien $S$. Soit $x_0 : \Spec(k) \to \mathcal{X}$
un morphisme, où $k$ est un corps de type fini sur $S$, dont
l'image est $s \in S$. Si $\mathcal{O}_{S, s}$ est un anneau G, alors
l'application du Lemme \ref{lemma-branches} préserve les multiplicités.
\end{lemma}

\begin{proof}
D'après le Lemme \ref{lemma-Artin-approximation-by-smooth-morphism}, nous
pouvons supposer qu'il existe un morphisme lisse $U \to \mathcal{X}$,
où $U$ est un schéma, et où un point à valeurs dans $k$, $u_0$, appartient à $U$,
tels que $\mathcal{O}_{U, u_0}^\wedge$ soit un anneau versel de
$\mathcal{X}$ en $x_0$. Par la construction de notre application dans
la preuve du Lemme \ref{lemma-branches} (qui se simplifie
considérablement puisque $A = \mathcal{O}_{U, u_0}^\wedge$), nous trouvons
qu'il suffit de montrer ceci : la multiplicité d'une composante irréductible
de $U$ passant par $u_0$ est égale à la
multiplicité de toute composante irréductible de
$\Spec(\mathcal{O}_{U, u_0}^\wedge)$ qui s'y envoie.

\medskip\noindent
Traduit en algèbre commutative, cela donne ce qui suit :
soit $C = \mathcal{O}_{U, u_0}$. Il est essentiellement de type fini
sur $\mathcal{O}_{S, s}$ et est donc un anneau G (Compléments d'algèbre,
Proposition \ref{more-algebra-proposition-finite-type-over-G-ring}).
Alors $A = C^\wedge$. Par conséquent, $C \to A$ est un homomorphisme régulier d'anneaux.
Soit $\mathfrak q \subset C$ un idéal premier minimal et soit $\mathfrak p \subset A$
un idéal premier minimal au-dessus de $\mathfrak q$. Alors
$$
R = C_\mathfrak p \longrightarrow A_\mathfrak p = R'
$$
est un homomorphisme régulier d'anneaux locaux artiniens. Pour un tel homomorphisme,
on a toujours
$$
\text{length}_R R = \text{length}_{R'} R'
$$
C'est précisément ce qu'il fallait montrer, car le membre de gauche
est la multiplicité de notre composante sur $U$, et le membre de droite
est la multiplicité de notre composante sur $\Spec(A)$.
Pour établir l'égalité, nous utilisons d'abord
$$
\text{length}_R(R) \text{length}_{R'}(R'/\mathfrak m_R R') =
\text{length}_{R'}(R')
$$
d'après Algèbre, Lemme \ref{algebra-lemma-pullback-module}. Il suffit donc
de montrer que $\mathfrak m_R R' = \mathfrak m_{R'}$, ce qui résulte
du fait qu'il s'agit d'un homomorphisme régulier d'anneaux locaux de dimension zéro.
\end{proof}













\section{Théorie de la dimension des champs algébriques}
\label{section-dimension-of-algebraic-stacks}

\noindent
À notre connaissance, les principaux résultats publiés sur la théorie de la
dimension des champs algébriques sont ceux de \cite{Osserman}, qui
étudie les notions de codimension et de dimension relative. Nous
examinons plus en détail la notion de dimension d'un
champ algébrique en un point, et démontrons divers résultats
reliant la dimension des fibres d'un morphisme en un point de la source
à la dimension de sa source et de son but. Nous démontrons aussi un résultat
(le Lemme \ref{lemma-dimension-formula} ci-dessous) qui
nous permet (sous des hypothèses appropriées) de calculer la dimension
d'un champ algébrique en un point au moyen d'un anneau versel.

\medskip\noindent
Sans avoir toujours cherché à optimiser nos résultats, nous avons
dans une large mesure essayé d'éviter les hypothèses superflues. Toutefois, dans certains
de nos résultats, où nous comparons certaines propriétés d'un champ
algébrique à celles d'un anneau versel de ce
champ en un point, nous nous sommes restreints
au cas des champs algébriques localement de présentation finie
sur un schéma de base localement noethérien dont tous les anneaux locaux sont
des anneaux $G$. Cela nous permet de disposer de l'approximation d'Artin
pour comparer la géométrie de l'anneau versel à celle
du champ lui-même. Il se peut cependant que cette hypothèse restrictive
ne soit pas nécessaire à la validité
de tous les énoncés que nous démontrons.
Comme elle est satisfaite dans les applications que nous avons en vue,
nous nous en sommes néanmoins accommodés lorsqu'elle est utile.

\medskip\noindent
Si $X$ est un schéma, nous définissons la dimension $\dim(X)$ de $X$
comme la dimension de Krull de l'
espace topologique sous-jacent à $X$,
tandis que, si $x$ est un point de $X$,
nous définissons la dimension $\dim_x (X)$ de $X$ en $x$ comme le
minimum des dimensions des sous-ensembles ouverts $U$ de $X$ contenant
$x$, voir
Propriétés, Définition \ref{properties-definition-dimension}.
On a la relation $\dim(X) = \sup_{x \in X} \dim_x(X)$, voir
Propriétés, Lemme \ref{properties-lemma-dimension}.
Si $X$ est localement noethérien, alors $\dim_x(X)$ coïncide avec le supremum
des dimensions en $x$
des composantes irréductibles de $X$ passant par $x$.

\medskip\noindent
Si $X$ est un espace algébrique et $x \in |X|$,
nous définissons $\dim_x X = \dim_u U,$ où $U$ est un schéma quelconque
admettant une surjection étale $U \to X$,
et où $u\in U$ est un point quelconque au-dessus de $x$, voir
Propriétés des espaces, Définition
\ref{spaces-properties-definition-dimension-at-point}.
Nous posons $\dim(X) = \sup_{x \in |X|} \dim_x(X)$, voir
Propriétés des espaces, Définition \ref{spaces-properties-definition-dimension}.

\begin{remark}
\label{remark-dimension-algebraic-space}
En général, la dimension de l'espace algébrique $X$ en un point $x$
peut ne pas coïncider avec la dimension de l'espace topologique sous-jacent
$|X|$ en $x$. Par exemple, si $k$ est un corps de caractéristique zéro et
$X =  \mathbf{A}^1_k / \mathbf{Z}$, alors $X$ est de dimension $1$ (la dimension
de $\mathbf{A}^1_k$) en chacun de ses points,
tandis que $|X|$ est muni de la topologie grossière et a donc une
dimension de Krull nulle. D'autre part, dans
Espaces algébriques, Exemple \ref{spaces-example-infinite-product},
on donne un exemple d'espace algébrique
qui est de dimension $0$ en chacun de ses points, tandis que $|X|$ est
irréductible, de dimension de Krull $1$, et admet un point générique (de sorte que la
dimension de $|X|$ en chacun de ses points est $1$); voir aussi la discussion
de cet exemple dans
Propriétés des espaces, section \ref{spaces-properties-section-dimension}.

\medskip\noindent
D'autre part, si $X$ est un espace algébrique {\it convenable}, au sens de
Espaces convenables, Définition \ref{decent-spaces-definition-very-reasonable}
(en particulier, si $X$ est quasi-séparé; voir
Espaces convenables, section \ref{decent-spaces-section-reasonable-decent}),
alors la dimension de $X$ en $x$ coïncide effectivement avec la dimension
de $|X|$ en $x$; voir
Espaces convenables, Lemme \ref{decent-spaces-lemma-dimension-decent-space}.
\end{remark}

\noindent
Pour définir la dimension d'un champ algébrique,
il sera utile de disposer d'abord de la notion de dimension relative,
en un point de la source,
d'un morphisme dont la source est un espace algébrique
et le but un champ algébrique. La définition est quelque peu
compliquée, simplement parce que (contrairement au cas des schémas) les points d'un
champ algébrique, ou d'un espace
algébrique, ne peuvent être décrits comme des morphismes depuis le spectre d'un corps,
mais seulement comme des classes d'équivalence de tels morphismes.

\begin{definition}
\label{definition-relative-dimension}
Si $f : T \to \mathcal{X}$ est un morphisme localement de type fini d'un
espace algébrique vers un champ algébrique,
et si $t \in |T|$ est un point d'image $x \in | \mathcal{X}|$, nous définissons
{\it la dimension relative} de $f$ en $t$, notée
$\dim_t(T_x),$
comme suit :
choisissons un morphisme $\Spec k \to \mathcal{X}$, de source le spectre d'un
corps, qui représente $x$, et choisissons un point
$t' \in |T \times_{\mathcal{X}} \Spec k|$
s'envoyant sur $t$ par la projection vers $|T|$
(un tel point $t'$ existe, d'après
Propriétés des champs, Lemme \ref{stacks-properties-lemma-points-cartesian});
alors
$$
\dim_t(T_x) = \dim_{t'}(T \times_{\mathcal{X}} \Spec k ).
$$
\end{definition}

\noindent
Notons que, puisque $T$ est un espace algébrique et $\mathcal{X}$ un
champ algébrique, le produit fibré $T \times_{\mathcal{X}} \Spec k$
est un espace algébrique; la quantité figurant au membre de droite de
cette définition est donc bien définie (voir la discussion ci-dessus).

\begin{remark}
\label{remark-relative-dimension}
(1)
On vérifie aisément (par exemple en utilisant l'invariance
par changement de base de la dimension relative des morphismes de schémas localement
de type fini; voir par exemple
Morphismes, Lemme \ref{morphisms-lemma-dimension-fibre-after-base-change})
que $\dim_t(T_x)$ est bien définie, indépendamment des choix
faits pour la calculer.

\medskip\noindent
(2)
Lorsque $\mathcal{X}$ est aussi un espace algébrique,
on vérifie immédiatement que cette définition coïncide avec
la définition de la dimension relative donnée dans
Morphismes d'espaces, Définition
\ref{spaces-morphisms-definition-dimension-fibre}.
\end{remark}

\noindent
Rappelons maintenant le lemme suivant, sur lequel repose notre étude de
la dimension d'un champ algébrique localement noethérien.

\begin{lemma}
\label{lemma-behaviour-of-dimensions-wrt-smooth-morphisms}
Si $f: U \to X$ est un morphisme lisse d'espaces algébriques localement
noethériens, et
si $u \in |U|$ est d'image $x \in |X|$, alors
$$
\dim_u (U) = \dim_x(X) + \dim_{u} (U_x)
$$
où $\dim_u (U_x)$ est défini par la
Définition \ref{definition-relative-dimension}.
\end{lemma}

\begin{proof}
Voir Morphismes d'espaces, Lemme
\ref{spaces-morphisms-lemma-smoothness-dimension-spaces}
en notant que la définition de $\dim_u (U_x)$ utilisée ici coïncide avec
celle qui y est employée, d'après la Remarque \ref{remark-relative-dimension} (2).
\end{proof}

\begin{lemma}
\label{lemma-dimension-for-stacks}
Soient $\mathcal{X}$ un champ algébrique localement noethérien et
$x \in |\mathcal{X}|$. Soit $U \to \mathcal{X}$ un morphisme lisse
d'un espace algébrique vers $\mathcal{X}$, et soit $u$ un point quelconque de $|U|$
s'envoyant sur $x$. Alors
$$
\dim_x(\mathcal{X}) =  \dim_u(U) - \dim_{u}(U_x)
$$
où la dimension relative $\dim_u(U_x)$ est définie
par la Définition \ref{definition-relative-dimension}, et la
dimension de $\mathcal{X}$ en $x$ comme dans
Propriétés des champs, Définition
\ref{stacks-properties-definition-dimension-at-point}.
\end{lemma}

\begin{proof}
Le Lemme \ref{lemma-behaviour-of-dimensions-wrt-smooth-morphisms}
permet de vérifier que le membre de droite
$\dim_u(U) + \dim_u(U_x)$ est indépendant du choix
du morphisme lisse $U \to \mathcal{X}$ et de $u \in |U|$.
Nous omettons les détails. En particulier, nous pouvons supposer que $U$ est un schéma.
Dans ce cas, nous pouvons calculer $\dim_u(U_x)$
en choisissant comme représentant de $x$
le composé $\Spec \kappa(u) \to U \to \mathcal{X}$, où
le premier morphisme est le morphisme canonique d'image $u \in U$.
Écrivons alors $R = U \times_{\mathcal{X}} U$ et notons
$e : U \to R$ le morphisme diagonal. L'invariance de la
dimension relative par changement de base montre que
$\dim_u(U_x) = \dim_{e(u)}(R_u)$. Nous voyons donc que
le membre de droite est égal à
$\dim_u (U) - \dim_{e(u)}(R_u) = \dim_x(\mathcal{X})$, comme voulu.
\end{proof}

\begin{remark}
\label{remark-dimension-DM}
Pour les champs de Deligne--Mumford suffisamment convenables
(par exemple quasi-séparés),
$\dim_x(\mathcal{X})$ coïncide à nouveau avec la
quantité définie
topologiquement $\dim_x |\mathcal{X}|$. Cependant, pour des champs d'Artin plus
généraux,
ce n'est généralement pas le cas. Par exemple, si
$\mathcal{X} = [\mathbf{A}^1/\mathbf{G}_m]$
(sur un certain corps, le quotient étant pris pour
l'action multiplicative usuelle de $\mathbf{G}_m$ sur $\mathbf{A}^1$),
alors $|\mathcal{X}|$ possède deux points, dont l'un est une spécialisation de l'autre
(correspondant
aux deux orbites de $\mathbf{G}_m$ sur $\mathbf{A}^1$), et est donc de
dimension $1$ comme
espace topologique; mais $\dim_x (\mathcal{X}) = 0$ pour les deux points
$x \in |\mathcal{X}|$.
(Un exemple encore plus extrême est fourni par l'espace classifiant
$[\Spec k/\mathbf{G}_m]$, dont la dimension en son unique point
est égale à $-1$.)
\end{remark}

\noindent
Nous pouvons maintenant étendre la Définition \ref{definition-relative-dimension}
au contexte des morphismes (localement de type fini)
entre champs algébriques (localement noethériens).

\begin{definition}
\label{definition-relative-dimension-for-stacks}
Si $f : \mathcal{T} \to \mathcal{X}$
est un morphisme localement de type fini entre
champs algébriques localement noethériens, et si
$t \in |\mathcal{T}|$ est un point d'image $x \in |\mathcal{X}|$, nous
définissons la {\it dimension relative} de $f$ en $t$, notée
$\dim_t(\mathcal{T}_x),$ comme suit :
choisissons un morphisme $\Spec k \to \mathcal{X}$, de source le spectre d'un
corps, qui représente $x$, et choisissons un point
$t' \in |\mathcal{T} \times_{\mathcal{X}} \Spec k|$
s'envoyant sur $t$ par la projection vers $|\mathcal{T}|$
(un tel point $t'$ existe, d'après
Propriétés des champs, Lemme
\ref{stacks-properties-lemma-points-cartesian}; alors
$$
\dim_t(\mathcal{T}_x) = \dim_{t'}(\mathcal{T} \times_{\mathcal{X}} \Spec k ).
$$
\end{definition}

\noindent
Notons que, puisque $\mathcal{T}$ est un champ algébrique et $\mathcal{X}$ un
champ algébrique, le produit fibré $\mathcal{T}\times_{\mathcal{X}} \Spec k$
est un champ algébrique, localement noethérien d'après
Morphismes de champs, Lemme
\ref{stacks-morphisms-lemma-locally-finite-type-locally-noetherian}.
Ainsi, la quantité figurant au membre de droite de cette définition
est définie par Propriétés des champs,
Définition \ref{stacks-properties-definition-dimension-at-point}.

\begin{remark}
\label{remark-dimension-tangent-space-well-defined}
Des manipulations usuelles montrent que $\dim_t(\mathcal{T}_x)$ est bien définie,
indépendamment des choix faits pour la calculer.
\end{remark}

\noindent
Établissons maintenant quelques propriétés élémentaires de la dimension relative,
qui sont des généralisations évidentes des énoncés correspondants dans le
cas des morphismes de schémas.

\begin{lemma}
\label{lemma-base-change-invariance-of-relative-dimension}
Supposons donné
un carré cartésien de morphismes de champs localement noethériens
$$
\xymatrix{
\mathcal{T}' \ar[d]\ar[r] & \mathcal{T} \ar[d] \\
\mathcal{X}' \ar[r] & \mathcal{X}
}
$$
où les morphismes verticaux sont localement de type fini.
Si $t' \in |\mathcal{T}'|$,
a pour images $t$, $x'$ et $x$ dans $|\mathcal{T}|$, $|\mathcal{X}'|$ et
$|\mathcal{X}|$,
respectivement, alors $\dim_{t'}(\mathcal{T}'_{x'}) = \dim_{t}(\mathcal{T}_x).$
\end{lemma}

\begin{proof}
Les deux membres peuvent (par définition) être calculés comme la
dimension du même produit fibré.
\end{proof}

\begin{lemma}
\label{lemma-behaviour-of-dimensions-wrt-smooth-morphisms-stacky}
Si $f: \mathcal{U} \to \mathcal{X}$ est un morphisme lisse de champs algébriques
localement noethériens, et
si $u \in |\mathcal{U}|$ est d'image $x \in |\mathcal{X}|$,
alors
$$
\dim_u (\mathcal{U}) = \dim_x(\mathcal{X}) + \dim_{u} (\mathcal{U}_x).
$$
\end{lemma}

\begin{proof}
Choisissons un morphisme lisse surjectif $V \to \mathcal{U}$ dont la source
est un schéma, et soit $v\in |V|$ un point s'envoyant sur $u$.
Alors le composé $V \to \mathcal{U} \to \mathcal{X}$ est lui aussi lisse,
et, d'après le Lemme \ref{lemma-behaviour-of-dimensions-wrt-smooth-morphisms},
nous avons $\dim_x(\mathcal{X}) = \dim_v(V) - \dim_v(V_x),$
tandis que $\dim_u(\mathcal{U}) = \dim_v(V) - \dim_v(V_u).$
Ainsi
$$
\dim_u(\mathcal{U}) - \dim_x(\mathcal{X}) = \dim_v (V_x) - \dim_v (V_u).
$$

\medskip\noindent
Choisissons un représentant $\Spec k \to \mathcal{X}$ de $x$
et un point $v' \in | V \times_{\mathcal{X}} \Spec k|$ au-dessus de
$v$, d'image $u'$ dans $|\mathcal{U}\times_{\mathcal{X}} \Spec k|$;
alors, par définition,
$\dim_u(\mathcal{U}_x) = \dim_{u'}(\mathcal{U}\times_{\mathcal{X}} \Spec k),$
et
$\dim_v(V_x) = \dim_{v'}(V\times_{\mathcal{X}} \Spec k).$

\medskip\noindent
Or $V\times_{\mathcal{X}} \Spec k \to \mathcal{U}\times_{\mathcal{X}}\Spec k$
est un morphisme lisse surjectif (comme changement de base
d'un tel morphisme) dont la source est un espace algébrique
(puisque $V$ et $\Spec k$ sont des schémas et $\mathcal{X}$
un champ algébrique). Ainsi, toujours par définition,
nous avons
\begin{align*}
\dim_{u'}(\mathcal{U}\times_{\mathcal{X}} \Spec k)
& =
\dim_{v'}(V\times_{\mathcal{X}} \Spec k) -
\dim_{v'}(V \times_{\mathcal{X}} \Spec k)_{u'}) \\
& = \dim_v(V_x) -
\dim_{v'}( (V\times_{\mathcal{X}} \Spec k)_{u'}).
\end{align*}
Or $V\times_{\mathcal{X}} \Spec k \cong
V\times_{\mathcal{U}} (\mathcal{U}\times_{\mathcal{X}} \Spec k),$
et donc le
Lemme \ref{lemma-base-change-invariance-of-relative-dimension}
montre que
$\dim_{v'}((V\times_{\mathcal{X}} \Spec k)_{u'})  = \dim_v(V_u).$
En réunissant ces égalités, nous trouvons
$$
\dim_u(\mathcal{U}) - \dim_x(\mathcal{X}) =
\dim_u(\mathcal{U}_x),
$$
comme voulu.
\end{proof}

\begin{lemma}
\label{lemma-relative-dimension-is-semi-continuous}
Soit $f: \mathcal{T} \to \mathcal{X}$ un morphisme localement de type fini de
champs algébriques.
\begin{enumerate}
\item
La fonction $t \mapsto \dim_t(\mathcal{T}_{f(t)})$ est semi-continue supérieurement
sur $|\mathcal{T}|$.
\item Si $f$ est lisse, alors
la fonction $t \mapsto \dim_t(\mathcal{T}_{f(t)})$ est localement constante
sur $|\mathcal{T}|$.
\end{enumerate}
\end{lemma}

\begin{proof}
Supposons d'abord que $\mathcal{T}$ soit un schéma $T$,
soit $U \to \mathcal{X}$ un morphisme lisse surjectif dont la source
est un schéma, et soit $T' = T \times_{\mathcal{X}} U$.
Soit $f': T' \to U$ le changement de base de $f$ à $U$,
et soit $g: T' \to T$ la projection.

\medskip\noindent
Le Lemme \ref{lemma-base-change-invariance-of-relative-dimension}
montre que $\dim_{t'}(T'_{f'(t')}) = \dim_{g(t')}(T_{f(g(t'))}),$
pour $t' \in T'$, tandis que,
puisque $g$ est lisse et surjectif (étant le changement de base
d'un morphisme lisse surjectif), l'application induite par $g$ sur les espaces
topologiques sous-jacents est continue et ouverte
(d'après
Propriétés des espaces, Lemme \ref{spaces-properties-lemma-topology-points}), et
surjective. Il suffit donc de noter que la partie (1) pour le morphisme $f'$
résulte de
Morphismes d'espaces, Lemme
\ref{spaces-morphisms-lemma-openness-bounded-dimension-fibres}, et la partie (2)
de l'un ou l'autre des lemmes suivants, Morphismes, Lemme
\ref{morphisms-lemma-flat-finite-presentation-CM-fibres-relative-dimension}
ou
Morphismes, Lemme \ref{morphisms-lemma-smooth-omega-finite-locally-free}
(chacun donne le résultat pour les schémas, d'où
les résultats analogues pour les espaces algébriques se
déduisent exactement comme dans
Morphismes d'espaces, Lemme
\ref{spaces-morphisms-lemma-openness-bounded-dimension-fibres}.

\medskip\noindent
Revenons maintenant au cas général,
et choisissons un morphisme lisse surjectif
$h:V \to \mathcal{T}$ dont la source est un schéma.
Si $v \in V$, alors, essentiellement par définition,
nous avons
$$
\dim_{h(v)}(\mathcal{T}_{f(h(v))}) =
\dim_{v}(V_{f(h(v))}) - \dim_{v}(V_{h(v)}).
$$
Puisque $V$ est un schéma, nous avons démontré que le premier
des termes du membre de droite de cette égalité
est semi-continu supérieur (et même localement
constant si $f$ est lisse), tandis que le second terme est
en fait localement constant.
Leur différence est donc semi-continue supérieurement
(et localement constante si $f$ est lisse),
et par conséquent la fonction
$\dim_{h(v)}(\mathcal{T}_{f(h(v))})$
est semi-continue supérieurement sur $|V|$ (et localement
constante si $f$ est lisse).
Puisque le morphisme $|V| \to |\mathcal{T}|$ est ouvert et surjectif,
le lemme en résulte.
\end{proof}

\noindent
Avant de poursuivre notre étude,
démontrons deux lemmes relatifs à la théorie de la dimension des schémas.

\medskip\noindent
Pour situer le premier lemme,
notons que, si $X$ est un schéma de dimension finie, puisque $\dim X$
est défini comme le supremum des dimensions $\dim_x X$,
il existe un point $x \in X$ tel que $\dim_x X = \dim X$.
Le lemme suivant montre que nous pouvons en outre choisir le point
$x$ de type fini.

\begin{lemma}
\label{lemma-dimension-achieved-by-finite-type-point}
Si $X$ est un schéma de dimension finie,
il existe un point fermé (et donc de type fini) $x \in X$
tel que $\dim_x X = \dim X$.
\end{lemma}

\begin{proof}
Soit $d = \dim X$,
et choisissons une chaîne strictement décroissante maximale
de sous-ensembles fermés irréductibles de $X$,
disons
\begin{equation}
\label{equation-maximal-chain}
Z_0 \supset Z_1 \supset \ldots \supset Z_d.
\end{equation}
Le sous-ensemble $Z_d$ est un sous-ensemble fermé irréductible minimal de $X$,
et tout point de $Z_d$ est donc un point générique de $Z_d$.
Puisque l'espace topologique sous-jacent au schéma $X$ est sobre,
nous concluons que $Z_d$ est un singleton, constitué d'un unique
point fermé $x \in X$.
Si $U$ est
un voisinage quelconque de $x$, alors
la chaîne
$$
U\cap Z_0 \supset U\cap Z_1 \supset \ldots \supset U\cap Z_d = Z_d =
\{x\}
$$
est une chaîne strictement décroissante de sous-ensembles
fermés irréductibles de $U$, ce qui montre que $\dim U \geq d$.
Nous trouvons donc $\dim_x X \geq d$. L'autre inégalité
étant évidente, le lemme est démontré.
\end{proof}

\noindent
Le lemme suivant montre que $\dim_x X$ est une fonction {\it constante}
sur un schéma irréductible satisfaisant quelques hypothèses supplémentaires modérées.

\begin{lemma}
\label{lemma-constancy-of-dimension}
Si $X$ est un schéma irréductible, jacobsonien, caténaire et localement
noethérien, de dimension finie,
alors $\dim U = \dim X$ pour tout
sous-ensemble ouvert non vide $U$ de $X$.
De manière équivalente, $\dim_x X$ est une fonction constante sur $X$.
\end{lemma}

\begin{proof}
L'équivalence des deux assertions résulte directement des
définitions. Supposons donc que $U\subset X$ soit un sous-ensemble ouvert
non vide.
Certes $\dim U \leq \dim X$, et nous devons montrer
que $\dim U \geq \dim X.$
Écrivons $d = \dim X$, et choisissons une chaîne strictement
décroissante maximale de sous-ensembles fermés irréductibles
de $X$, disons
$$
X = Z_0 \supset Z_1 \supset \ldots \supset Z_d.
$$
Puisque $X$ est jacobsonien, le sous-ensemble fermé irréductible
minimal $Z_d$ est égal à $\{x\}$ pour un certain
point fermé $x$.

\medskip\noindent
Si $x \in U,$ alors
$$
U = U \cap Z_0  \supset U\cap Z_1 \supset \ldots \supset
U\cap Z_d = \{x\}
$$
est une chaîne strictement décroissante de sous-ensembles fermés
irréductibles de $U$, et nous concluons donc que $\dim U \geq d$,
comme voulu. Nous pouvons donc supposer que $x \not\in U.$

\medskip\noindent
Considérons le morphisme plat $\Spec \mathcal{O}_{X,x} \to X$.
Le sous-ensemble ouvert non vide (et donc dense) $U$ de $X$
se relève en un sous-ensemble ouvert $V \subset \Spec \mathcal{O}_{X,x}$.
En remplaçant $U$ par un sous-ensemble ouvert non vide quasi-compact, donc
noethérien, nous pouvons supposer que l'inclusion
$U \to X$ est un morphisme quasi-compact. Puisque la
formation des images schématiques de morphismes quasi-compacts
commute au changement de base plat, d'après
Morphismes, Lemme
\ref{morphisms-lemma-flat-base-change-scheme-theoretic-image}
nous voyons que $V$ est dense dans $\Spec \mathcal{O}_{X,x}$,
et en particulier non vide,
et bien sûr $x \not\in V.$ (Ici, nous notons aussi $x$
le point fermé de $\Spec \mathcal{O}_{X,x}$, puisque son image
est égale au point donné $x \in X$.)
Or $\Spec \mathcal{O}_{X,x} \setminus \{x\}$ est jacobsonien d'après
Propriétés, Lemme
\ref{properties-lemma-complement-closed-point-Jacobson}
et $V$ contient donc un point fermé $z$
de $\Spec \mathcal{O}_{X,x} \setminus \{x\}$. L'adhérence
dans $X$ de l'image de $z$ est alors un sous-ensemble
fermé irréductible $Z$ de $X$ contenant $x$, dont l'intersection
avec $U$ est non vide, et
pour lequel il n'existe aucun sous-ensemble fermé
irréductible strictement contenu dans $Z$
et contenant strictement $\{x\}$
(car le changement de base à $\Spec \mathcal{O}_{X,x}$ induit
une bijection entre les sous-ensembles fermés irréductibles de $X$
contenant $x$ et les sous-ensembles fermés irréductibles de $\Spec
\mathcal{O}_{X,x}$).
Puisque $U \cap Z$ est un sous-ensemble fermé non vide de $U$,
il contient un point $u$ qui est fermé dans $X$ (puisque
$X$ est jacobsonien), et puisque $U\cap Z$
est un sous-ensemble ouvert non vide (donc dense) de l'ensemble irréductible $Z$
(qui contient un point n'appartenant pas à $U$, à savoir $x$),
l'inclusion $\{u\} \subset U\cap Z$ est stricte.

\medskip\noindent
Puisque $X$ est caténaire, la chaîne
$$
X = Z_0 \supset Z \supset \{x\} = Z_d
$$
peut être raffinée en une chaîne de longueur $d+1$, qui est nécessairement
de la forme
$$
X = Z_0 \supset W_1 \supset \ldots \supset W_{d-1} = Z \supset \{x\} = Z_d.
$$
Puisque $U\cap Z$ est non vide, nous trouvons alors que
$$
U = U \cap Z_0 \supset U \cap W_1\supset \ldots \supset U\cap W_{d-1}
= U\cap Z \supset \{u\}
$$
est une chaîne strictement décroissante de sous-ensembles fermés irréductibles
de $U$, de longueur $d+1$, ce qui montre que $\dim U \geq d$,
comme voulu.
\end{proof}

\noindent
Nous démontrerons un analogue pour les champs
du Lemme \ref{lemma-constancy-of-dimension}
dans le Lemme \ref{lemma-irreducible-implies-equidimensional} ci-dessous,
mais nous devons auparavant introduire une définition supplémentaire,
rendue nécessaire par le fait que la propriété d'un schéma d'être caténaire
n'est pas locale pour la topologie étale
(voir l'exemple de
Algèbre, Remarque \ref{algebra-remark-universally-catenary-does-not-descend},
ce qui rend difficile de définir ce que signifie, pour un espace algébrique
ou un champ algébrique, d'être caténaire
(voir la discussion de \cite[page 3]{Osserman}).
Pour certains aspects de la théorie de la dimension, la définition
suivante semble fournir un bon substitut à la notion manquante
de champ algébrique caténaire.

\begin{definition}
\label{definition-pseudo-catenary}
Nous disons qu'un champ algébrique localement noethérien $\mathcal{X}$
est {\it pseudo-caténaire} s'il existe un morphisme lisse
et surjectif $U \to \mathcal{X}$ dont la source est
un schéma universellement caténaire.
\end{definition}

\begin{example}
\label{example-pseudo-catenary}
Si $\mathcal{X}$ est localement de type fini sur un schéma universellement
caténaire et localement noethérien $S$,
et si $U\to \mathcal{X}$ est un morphisme lisse surjectif
dont la source est un schéma, alors le composé
$U \to \mathcal{X} \to S$ est localement de type fini,
et donc $U$ est universellement caténaire d'après
Morphismes, Lemme
\ref{morphisms-lemma-universally-catenary-local}.
Ainsi, $\mathcal{X}$ est pseudo-caténaire.
\end{example}

\noindent
Le lemme suivant montre que la propriété d'être pseudo-caténaire
se transmet par les morphismes de type fini.

\begin{lemma}
\label{lemma-catenary-covers}
Si $\mathcal{X}$ est un champ algébrique localement noethérien
pseudo-caténaire, et si $\mathcal{Y} \to \mathcal{X}$ est un morphisme
localement de type fini,
alors il existe un morphisme lisse surjectif $V \to \mathcal{Y}$
dont la source est un schéma universellement caténaire; ainsi
$\mathcal{Y}$ est encore pseudo-caténaire.
\end{lemma}

\begin{proof}
Par hypothèse, nous pouvons trouver un morphisme lisse surjectif
$U \to \mathcal{X}$ dont la source est un schéma universellement caténaire.
Le changement de base $U\times_{\mathcal{X}} \mathcal{Y}$ est alors un champ
algébrique; soit $V \to U\times_{\mathcal{X}} \mathcal{Y}$ un morphisme lisse
surjectif dont la source est un schéma.
Le composé $V \to U\times_{\mathcal{X}} \mathcal{Y} \to \mathcal{Y}$ est alors
lisse et surjectif (comme composé de morphismes lisses et
surjectifs), tandis que le morphisme $V \to U\times_{\mathcal{X}}
\mathcal{Y} \to U$ est localement de type fini (comme composé
de morphismes localement de type fini). Puisque $U$
est universellement caténaire, nous voyons que $V$ l'est aussi
(d'après Morphismes, Lemme
\ref{morphisms-lemma-universally-catenary-local}),
comme affirmé.
\end{proof}

\noindent
Étudions maintenant le comportement de la fonction $\dim_x(\mathcal{X})$ sur
$|\mathcal{X}|$
(pour un champ localement noethérien $\mathcal{X}$) par rapport aux
composantes
irréductibles de $|\mathcal{X}|$, ainsi que diverses
questions connexes.

\begin{lemma}
\label{lemma-irreducible-implies-equidimensional}
Si $\mathcal{X}$ est
un champ algébrique jacobsonien, pseudo-caténaire et localement noethérien
tel que $|\mathcal{X}|$ soit irréductible,
alors $\dim_x(\mathcal{X})$ est une fonction constante sur $|\mathcal{X}|$.
\end{lemma}

\begin{proof}
Il suffit de montrer que $\dim_x(\mathcal{X})$ est localement constante sur
$|\mathcal{X}|$,
car elle sera alors nécessairement constante (puisque $|\mathcal{X}|$ est connexe,
étant irréductible). Puisque $\mathcal{X}$ est pseudo-caténaire,
nous pouvons trouver un morphisme lisse surjectif $U \to \mathcal{X}$, avec $U$
un schéma universellement caténaire. Si $\{U_i\}$ est un
recouvrement de $U$ par des sous-schémas ouverts quasi-compacts, nous pouvons remplacer
$U$ par $\coprod U_i,$, et
il suffit de montrer que
la fonction $u \mapsto \dim_{f(u)}(\mathcal{X})$ est localement constante sur $U_i$.
Puisque nous vérifions ceci pour un seul $U_i$ à la fois, omettons désormais l'indice,
et écrivons simplement $U$ plutôt que $U_i$.
Puisque $U$ est quasi-compact, il
est la réunion d'un nombre fini de composantes irréductibles,
disons $T_1 \cup \ldots \cup T_n$. Notons que chaque $T_i$ est jacobsonien,
caténaire et localement noethérien,
car c'est un sous-schéma fermé du schéma jacobsonien, caténaire et localement noethérien
$U$.

\medskip\noindent
D'après le Lemme \ref{lemma-behaviour-of-dimensions-wrt-smooth-morphisms}, nous avons
$\dim_{f(u)}(\mathcal{X}) = \dim_{u}(U) - \dim_{u}(U_{f(u)}).$
Le Lemme \ref{lemma-relative-dimension-is-semi-continuous} (2)
montre que le second terme du membre de droite est localement
constant sur $U$, puisque $f$ est lisse;
nous devons donc montrer que $\dim_u(U)$
est localement constante sur $U$. Puisque $\dim_u(U)$ est le maximum
des dimensions $\dim_u T_i$, lorsque $T_i$ parcourt les composantes
de $U$ contenant $u$, il suffit de montrer
que, si un point $u$ appartient à deux composantes distinctes,
disons $T_i$ et $T_j$ (avec $i \neq j$),
alors $\dim_u T_i = \dim_u T_j$,
puis de noter que $t\mapsto \dim_t T$ est une fonction constante
sur un schéma irréductible et jacobsonien,
caténaire et localement noethérien, noté $T$,
(comme il résulte du Lemme \ref{lemma-constancy-of-dimension}).

\medskip\noindent
Soient $V = T_i \setminus (\bigcup_{i' \neq i} T_{i'})$
et $W = T_j \setminus (\bigcup_{i' \neq j} T_{i'})$.
Chacun de $V$ et $W$ est alors un sous-ensemble ouvert non vide de $U$,
et a donc une image ouverte non vide dans $|\mathcal{X}|$. Puisque $|\mathcal{X}|$ est
irréductible,
ces deux sous-ensembles ouverts non vides de $|\mathcal{X}|$ ont une intersection
non vide.
Soit $x$ un point de cette intersection, et soient $v \in V$ et
$w\in W$ des points s'envoyant sur $x$.
Nous trouvons alors
$$
\dim T_i = \dim V = \dim_v (U) = \dim_x (\mathcal{X}) + \dim_v (U_x)
$$
et, de même,
$$
\dim T_j = \dim W = \dim_w (U) = \dim_x (\mathcal{X}) + \dim_w (U_x).
$$
Puisque $u \mapsto \dim_u (U_{f(u)})$ est localement constante sur $U$,
et puisque $T_i \cup T_j$ est connexe (comme réunion de deux ensembles irréductibles,
donc connexes, d'intersection non vide),
nous voyons que $\dim_v (U_x) = \dim_w(U_x)$,
et donc, en comparant les deux équations précédentes,
que $\dim T_i = \dim T_j$, comme voulu.
\end{proof}

\begin{lemma}
\label{lemma-closed-immersions}
Si $\mathcal{Z} \hookrightarrow \mathcal{X}$ est une immersion fermée
de champs algébriques localement noethériens,
et si $z \in |\mathcal{Z}|$ a pour image $x \in |\mathcal{X}|$,
alors $\dim_z (\mathcal{Z}) \leq \dim_x(\mathcal{X})$.
\end{lemma}

\begin{proof}
Choisissons un morphisme lisse surjectif
$U\to \mathcal{X}$ dont la source est un schéma;
le morphisme obtenu par changement de base
$V = U\times_{\mathcal{X}} \mathcal{Z} \to \mathcal{Z}$
est alors lui aussi lisse et surjectif, et la projection
$V \to U$ est une immersion fermée.
Si $v \in |V|$ s'envoie sur $z \in |\mathcal{Z}|$, et
si nous notons $u$ l'image de $v$ dans $|U|$,
alors, manifestement,
$\dim_v(V) \leq \dim_u(U)$,
tandis que
$\dim_v (V_z) = \dim_u(U_x)$,
d'après le Lemme \ref{lemma-base-change-invariance-of-relative-dimension}.
Ainsi,
$$
\dim_z(\mathcal{Z})  = \dim_v(V) - \dim_v(V_z)
\leq \dim_u(U) - \dim_u(U_x) = \dim_x(\mathcal{X}),
$$
comme affirmé.
\end{proof}

\begin{lemma}
\label{lemma-dimension-via-components}
Si $\mathcal{X}$ est un champ algébrique localement noethérien, et si
$x \in |\mathcal{X}|$,
alors $\dim_x(\mathcal{X}) = \sup_{\mathcal{T}} \{ \dim_x(\mathcal{T}) \} $,
où $\mathcal{T}$ parcourt toutes les composantes irréductibles
de $|\mathcal{X}|$ passant par $x$ (munies de leur
structure réduite induite).
\end{lemma}

\begin{proof}
Le Lemme \ref{lemma-closed-immersions}
montre que
$\dim_x (\mathcal{T}) \leq \dim_x(\mathcal{X})$ pour chaque
pour chaque composante irréductible $\mathcal{T}$ passant par
le point $x$. Pour démontrer le lemme,
il suffit donc de montrer que
\begin{equation}
\label{equation-desired-inequality}
\dim_x(\mathcal{X}) \leq
\sup_{\mathcal{T}} \{\dim_x(\mathcal{T})\}.
\end{equation}
Soit $U\to\mathcal{X}$ un recouvrement lisse par un schéma. Si $T$ est une composante
irréductible de $U$, notons $\mathcal{T}$ l'adhérence de son image
dans $\mathcal{X}$; c'est une composante irréductible de $\mathcal{X}$. Soit
$u \in U$
un point s'envoyant sur $x$. Alors
$\dim_x(\mathcal{X})=\dim_uU-\dim_uU_x=\sup_T\dim_uT-\dim_uU_x$, où le
supremum est pris sur les composantes irréductibles de $U$ passant
par $u$. Choisissons une composante $T$ sur laquelle le supremum
est atteint, et notons que
$\dim_x(\mathcal{T})=\dim_uT-\dim_u T_x$.
L'inégalité voulue (\ref{equation-desired-inequality})
résulte maintenant de l'inégalité évidente $\dim_u T_x \leq \dim_u U_x.$
(Notons que, si $\Spec k \to \mathcal{X}$ est un représentant de $x$,
alors $T\times_{\mathcal{X}} \Spec k$ est un sous-espace fermé de
$U\times_{\mathcal{X}}
\Spec k$.)
\end{proof}

\begin{lemma}
\label{lemma-dimension-at-finite-type-point}
Si $\mathcal{X}$ est un champ algébrique localement noethérien, et si
$x \in |\mathcal{X}|$, alors,
pour tout sous-champ ouvert $\mathcal{V}$ de $\mathcal{X}$ contenant $x$,
il existe un point de type fini $x_0 \in |\mathcal{V}|$ tel que
$\dim_{x_0}(\mathcal{X}) = \dim_x(\mathcal{V})$.
\end{lemma}

\begin{proof}
Choisissons un morphisme lisse surjectif
$f:U \to \mathcal{X}$ dont la source est un schéma, et considérons la
fonction $u \mapsto \dim_{f(u)}(\mathcal{X});$
puisque le morphisme $|U| \to |\mathcal{X}|$ induit par $f$ est ouvert (car $f$
est lisse) et surjectif (par hypothèse),
et envoie les points de type fini sur des points de type fini (par la définition même
des points de type fini de $|\mathcal{X}|$),
il suffit de montrer que, pour tout $u \in U$ et tout voisinage ouvert de $u$,
il existe dans ce voisinage un point de type fini $u_0$ tel que
$\dim_{f(u_0)}(\mathcal{X}) = \dim_{f(u)}(\mathcal{X}).$
Avec cette reformulation
du problème, la surjectivité de $f$ n'étant plus requise,
nous pouvons remplacer $U$ par le voisinage ouvert du point $u$ considéré,
et nous ramener ainsi à montrer que, pour chaque $u \in U$,
il existe un point de type fini $u_0 \in U$ tel que
$\dim_{f(u_0)}(\mathcal{X}) = \dim_{f(u)}(\mathcal{X}).$
D'après le Lemme \ref{lemma-behaviour-of-dimensions-wrt-smooth-morphisms},
$\dim_{f(u)}(\mathcal{X}) = \dim_u(U) - \dim_u(U_{f(u)}),$
tandis que
$\dim_{f(u_0)}(\mathcal{X}) = \dim_{u_0}(U) - \dim_{u_0}(U_{f(u_0)}).$
Puisque $f$ est lisse, l'expression $\dim_{u_0}(U_{f(u_0)})$ est localement
constante lorsque $u_0$ varie dans $U$ (d'après le
Lemme \ref{lemma-relative-dimension-is-semi-continuous} (2));
quitte donc à rétrécir encore $U$ autour de
$u$, nous pouvons la supposer constante. Le problème
se réduit ainsi à montrer que nous pouvons trouver un point de type fini $u_0 \in U$
pour lequel $\dim_{u_0}(U) = \dim_u(U)$.
Comme, par définition, $\dim_u U$ est le minimum des dimensions
$\dim V$, lorsque $V$ parcourt les voisinages ouverts $V$ de $u$
dans $U$, nous pouvons encore rétrécir $U$ autour de $u$ de sorte que
$\dim_u U = \dim U$.
L'existence du point voulu $u_0$ résulte alors du
Lemme \ref{lemma-dimension-achieved-by-finite-type-point}.
\end{proof}

\begin{lemma}
\label{lemma-monomorphing-a-component-in-of-the-right-dimension}
Soit $\mathcal{T} \hookrightarrow \mathcal{X}$ un monomorphisme localement
de type fini de champs algébriques,
où $\mathcal{X}$ (et donc aussi $\mathcal{T}$)
est jacobsonien, pseudo-caténaire et localement noethérien.
Supposons en outre que $\mathcal{T}$ soit irréductible,
de dimension (finie) $d$, et que $\mathcal{X}$ soit réduit
et de dimension inférieure
ou égale à $d$.
Alors il existe un sous-champ ouvert non vide $\mathcal{V}$ de $\mathcal{T}$ tel
que le monomorphisme
induit $\mathcal{V} \hookrightarrow \mathcal{X}$ soit une immersion ouverte
qui identifie
$\mathcal{V}$ à un sous-ensemble ouvert d'une composante irréductible de $\mathcal{X}$.
\end{lemma}

\begin{proof}
Choisissons un morphisme lisse surjectif $f:U \to \mathcal{X}$ de source un schéma,
nécessairement réduit puisque $\mathcal{X}$ l'est,
et écrivons $U' = \mathcal{T}\times_{\mathcal{X}} U$. Le morphisme obtenu par changement de base
$U' \to U$ est un monomorphisme d'espaces algébriques, localement de type
fini, donc représentable d'après
Morphismes d'espaces, Lemmes
\ref{spaces-morphisms-lemma-locally-quasi-finite-separated-representable} et
\ref{spaces-morphisms-lemma-monomorphism-loc-finite-type-loc-quasi-finite};
puisque $U$ est un schéma, $U'$ l'est aussi.
La projection $f': U' \to \mathcal{T}$ est encore une surjection lisse.
Soit $u' \in U'$, d'image $u \in U$.
Le Lemme \ref{lemma-base-change-invariance-of-relative-dimension}
montre que $\dim_{u'}(U'_{f(u')}) = \dim_u(U_{f(u)}),$
tandis que $\dim_{f'(u')}(\mathcal{T}) =d
\geq \dim_{f(u)}(\mathcal{X})$ d'après le
Lemme \ref{lemma-irreducible-implies-equidimensional}
et nos hypothèses sur $\mathcal{T}$ et $\mathcal{X}$.
Nous voyons donc que
\begin{equation}
\label{equation-dim-inequality}
\dim_{u'} (U') = \dim_{u'} (U'_{f(u')}) + \dim_{f'(u')}(\mathcal{T})
\\
\geq \dim_u (U_{f(u)}) + \dim_{f(u)}(\mathcal{X}) = \dim_u (U).
\end{equation}
Puisque $U' \to U$ est un monomorphisme localement de type fini,
il est en particulier non ramifié;
par la structure locale étale des morphismes non ramifiés,
Morphismes étales, Lemme \ref{etale-lemma-finite-unramified-etale-local},
nous pouvons donc trouver un diagramme commutatif
$$
\xymatrix{
V' \ar[r]\ar[d] & V \ar[d] \\
U' \ar[r] & U
}
$$
où le schéma $V'$ est non vide,
les flèches verticales sont étales,
et la flèche horizontale supérieure est une immersion fermée.
En remplaçant $V$ par un sous-ensemble ouvert quasi-compact
dont l'image rencontre l'image de $U'$,
et $V'$ par l'image réciproque de $V$, nous pouvons en outre
supposer que $V$ (et donc $V'$) est quasi-compact.
Puisque $V$ est aussi localement noethérien,
il est donc noethérien, et est par conséquent la réunion d'un nombre fini de composantes
irréductibles.

\medskip\noindent
Puisque les morphismes étales préservent la dimension ponctuelle, d'après
Descente, Lemme \ref{descent-lemma-dimension-at-point-local},
nous déduisons de (\ref{equation-dim-inequality})
que, pour tout point $v' \in V'$,
d'image $v \in V$, nous avons
$\dim_{v'}( V') \geq \dim_v(V)$.
En particulier, l'image de $V'$ ne peut être contenue dans l'intersection
de deux composantes irréductibles distinctes de $V$; nous pouvons donc trouver
au moins un sous-ensemble ouvert irréductible de $V$ ayant une intersection non vide
avec $V'$; en remplaçant $V$ par cet ouvert, nous pouvons supposer $V$ intègre
(car à la fois réduit et irréductible). L'inégalité précédente
sur les dimensions permet de conclure que l'immersion fermée $V' \hookrightarrow V$
est en fait un isomorphisme.
Si nous notons $W$ l'image de $V'$
dans $U'$, alors $W$ est un sous-ensemble
ouvert non vide de $U'$ (car les morphismes étales sont ouverts),
et le monomorphisme induit $W \to U$ est étale
(puisqu'il l'est localement pour la topologie étale sur la source, c'est-à-dire après changement de base
à $V'$), et est donc une immersion ouverte (étant un monomorphisme étale).
Ainsi, si nous notons $\mathcal{V}$ l'image de $W$ dans $\mathcal{T}$,
alors $\mathcal{V}$ est un sous-champ ouvert dense (de manière équivalente, non vide) de
$\mathcal{T}$,
dont l'image est dense dans une composante irréductible de $\mathcal{X}$.
Enfin,
notons que le morphisme $\mathcal{V} \to \mathcal{X}$ est lisse
(puisque son composé
avec le morphisme lisse $W\to \mathcal{V}$ est lisse),
et est aussi un monomorphisme, donc une immersion ouverte.
\end{proof}

\begin{lemma}
\label{lemma-dims-of-images}
Soit $f: \mathcal{T} \to \mathcal{X}$ un morphisme localement de type fini
de champs algébriques jacobsoniens, pseudo-caténaires et localement
noethériens,
dont la source est irréductible et le but quasi-séparé,
et soit $\mathcal{Z} \hookrightarrow \mathcal{X}$ l'image
schématique de $\mathcal{T}$.
Alors, pour tout $t \in |T|$,
nous avons
$\dim_t( \mathcal{T}_{f(t)}) \geq \dim \mathcal{T}  - \dim \mathcal{Z}$,
et il existe un sous-ensemble ouvert non vide (de manière équivalente, dense)
de $|\mathcal{T}|$ sur lequel l'égalité est satisfaite.
\end{lemma}

\begin{proof}
En remplaçant $\mathcal{X}$ par $\mathcal{Z}$, nous pouvons supposer, et supposons, que $f$ est
schématiquement dominant,
et aussi que $\mathcal{X}$ est irréductible.
Par la semi-continuité supérieure des dimensions des fibres
(Lemme \ref{lemma-relative-dimension-is-semi-continuous} (1)),
il suffit de démontrer que l'égalité
$\dim_t( \mathcal{T}_{f(t)}) =\dim \mathcal{T}  - \dim \mathcal{Z}$
est satisfaite pour $t$ appartenant à
un certain sous-champ ouvert non vide de $\mathcal{T}$.
Pour cette raison, nous pouvons à tout moment dans le raisonnement
remplacer $\mathcal{T}$ par un sous-champ ouvert non vide.

\medskip\noindent
Soit $T' \to \mathcal{T}$ un morphisme lisse surjectif dont la source
est un schéma, et soit $T$ un sous-ensemble ouvert non vide quasi-compact
de $T'$. Puisque $\mathcal{Y}$ est quasi-séparé, nous trouvons
que $T \to  \mathcal{Y}$ est quasi-compact
(d'après Morphismes de champs, Lemme
\ref{stacks-morphisms-lemma-quasi-compact-permanence}, appliqué aux morphismes
 $T \to \mathcal{Y} \to \Spec \mathbf{Z}$ auquel on l'applique).
Ainsi, en remplaçant $\mathcal{T}$ par l'image de $T$ dans $\mathcal{T}$,
nous pouvons supposer (en invoquant
Morphismes de champs, Lemme
\ref{stacks-morphisms-lemma-surjection-from-quasi-compact}
que le morphisme $f:\mathcal{T} \to \mathcal{X}$ est quasi-compact.

\medskip\noindent
Si nous choisissons une surjection lisse $U \to \mathcal{X}$ avec $U$ un schéma,
alors le Lemme \ref{lemma-map-of-components} assure que
nous pouvons trouver un sous-ensemble ouvert irréductible $V$ de $U$ tel
que $V \to \mathcal{X}$ soit lisse et schématiquement dominant.
Puisque la dominance schématique des morphismes quasi-compacts
est préservée par changement de base plat,
le changement de base $\mathcal{T} \times_{\mathcal{X}} V \to V$
du morphisme schématiquement
dominant $f$ est encore
schématiquement dominant. Soit $Z$ un schéma
admettant une surjection lisse sur ce produit fibré;
alors $Z \to \mathcal{T} \times_{\mathcal{X}} V \to V$
est encore schématiquement dominant.
Nous pouvons donc trouver une composante
irréductible $C$ de $Z$ qui domine
schématiquement $V$.
Puisque le composé $Z \to \mathcal{T}\times_{\mathcal{X}} V \to \mathcal{T}$
est lisse,
et puisque $\mathcal{T}$ est irréductible,
le Lemme \ref{lemma-map-of-components} montre que toute composante irréductible
de la source a une image dense dans $|\mathcal{T}|$.
Remplaçons maintenant
$C$ par un sous-ensemble ouvert non vide $W$ disjoint de toute autre
composante irréductible de $Z$, puis
remplaçons $\mathcal{T}$ et $\mathcal{X}$ par les images de $W$
et de $V$
(et appliquons le Lemme \ref{lemma-irreducible-implies-equidimensional}
pour voir que cela
ne change la dimension ni de $\mathcal{T}$ ni de $\mathcal{X}$).
Si nous notons $\mathcal{W}$ l'image du morphisme
$W \to \mathcal{T}\times_{\mathcal{X}} V$,
alors $\mathcal{W}$ est ouvert dans $\mathcal{T}\times_{\mathcal{X}} V$ (puisque le
morphisme $W \to \mathcal{T}\times_{\mathcal{X}} V$ est lisse),
et est irréductible (comme image d'un schéma
irréductible). Nous aboutissons ainsi à un diagramme commutatif
$$
\xymatrix{
W \ar[dr] \ar[r]  & \mathcal{W} \ar[r] \ar[d] & V \ar[d] \\
& \mathcal{T} \ar[r] & \mathcal{X}
}
$$
où $W$ et $V$ sont des schémas,
les flèches verticales sont lisses et surjectives,
les flèches diagonales et la flèche horizontale
supérieure de gauche sont lisses,
et le morphisme induit $\mathcal{W} \to \mathcal{T}\times_{\mathcal{X}} V$ est
une immersion ouverte. À l'aide de ce diagramme et des définitions
des diverses dimensions intervenant dans
l'énoncé du lemme, nous ramènerons la vérification
du lemme au cas des schémas, où elle est connue.

\medskip\noindent
Fixons $w \in |W|$, d'image $w' \in |\mathcal{W}|$,
d'image $t \in |\mathcal{T}|$, d'image $v$ dans $|V|$,
et d'image $x$ dans $|\mathcal{X}|$.
Essentiellement par définition (en utilisant le
fait que $\mathcal{W}$ est ouvert dans $\mathcal{T}\times_{\mathcal{X}} V$, et que
la fibre d'un changement de base est le changement de base de la fibre),
nous obtenons les égalités
$$
\dim_v V_x = \dim_{w'} \mathcal{W}_t
$$
et
$$
\dim_t \mathcal{T}_x = \dim_{w'} \mathcal{W}_v.
$$
D'après le Lemme \ref{lemma-behaviour-of-dimensions-wrt-smooth-morphisms}
(la flèche diagonale et la flèche verticale de droite
de notre diagramme réalisent $W$ et $V$ comme des recouvrements lisses par des
schémas des champs $\mathcal{T}$ et $\mathcal{X}$), nous trouvons
$$
\dim_t \mathcal{T} = \dim_w W - \dim_w W_t
$$
et
$$
\dim_x \mathcal{X} = \dim_v V - \dim_v V_x.
$$
En combinant ces égalités, nous trouvons
$$
\dim_t \mathcal{T}_x - \dim_t \mathcal{T} + \dim_x \mathcal{X}
= \dim_{w'} \mathcal{W}_v - \dim_w W + \dim_w W_t + \dim_v V -
\dim_{w'} \mathcal{W}_t
$$
Puisque $W \to \mathcal{W}$ est une surjection lisse, il en va de même
après changement de base par le morphisme $\Spec \kappa(v) \to V$
(en considérant $W \to \mathcal{W}$ comme un morphisme sur $V$),
et ce morphisme lisse donne la première des deux
égalités suivantes
$$
\dim_w W_v - \dim_{w'} \mathcal{W}_v = \dim_w (W_v)_{w'} = \dim_w W_{w'};
$$
la seconde égalité résulte d'une comparaison directe des
deux fibres en jeu.
De même, si nous considérons $W \to \mathcal{W}$ comme un morphisme de schémas
sur $\mathcal{T}$, et effectuons le changement de base par un représentant du point
$t \in |\mathcal{T}|$, nous obtenons les égalités
$$
\dim_w W_t - \dim_{w'} \mathcal{W}_t = \dim_w (W_t)_{w'} = \dim_w W_{w'}.
$$
En réunissant le tout, nous trouvons
$$
\dim_t \mathcal{T}_x - \dim_t \mathcal{T} + \dim_x \mathcal{X}
=  \dim_w W_v - \dim_w W + \dim_v V.
$$
Notre but est de montrer que le membre de gauche de cette égalité
s'annule sur un sous-ensemble ouvert non vide
de $t$. Lorsque $w$ parcourt un sous-ensemble ouvert non vide de $W$,
son image $t \in |\mathcal{T}|$ parcourt un sous-ensemble ouvert non vide
de $|\mathcal{T}|$ (puisque $W \to \mathcal{T}$ est lisse).

\medskip\noindent
Nous sommes donc ramenés à montrer que, si $W\to V$ est un
morphisme schématiquement dominant de schémas irréductibles localement
noethériens, localement de type fini,
alors il existe un sous-ensemble ouvert non vide de
points $w\in W$ tels que $\dim_w W_v =\dim_w W - \dim_v V$
(où $v$ désigne l'image de $w$ dans $V$).
C'est un fait classique,
dont nous rappelons la preuve pour la commodité du lecteur.

\medskip\noindent
Nous pouvons remplacer $W$ et $V$ par leurs sous-schémas réduits sous-jacents
sans modifier la validité (ou non) de cette équation,
et pouvons donc les supposer en fait intègres.
Puisque $\dim_w W_v$ est localement constante sur $W,$ quitte à remplacer $W$
par un sous-ensemble ouvert non vide, nous pouvons supposer que $\dim_w W_v$
est constante, disons égale à $d$. En choisissant cet ouvert affine,
nous pouvons aussi supposer que le morphisme $W\to V$ est de type fini.
Quitte à remplacer $V$ par un sous-ensemble ouvert non vide
(et à changer ensuite la base de $W$ à cet ouvert; le changement de base obtenu
est non vide, puisque le changement de base plat d'un morphisme quasi-compact
et schématiquement
dominant demeure schématiquement dominant),
nous pouvons en outre supposer que $W$ est plat sur $V$.
Le morphisme $W\to V$ est donc de dimension relative $d$
au sens de
Morphismes, Définition
\ref{morphisms-definition-relative-dimension-d}
et il résulte de
Morphismes, Lemme \ref{morphisms-lemma-rel-dimension-dimension},
que $\dim_w(W) = \dim_v(V) + d,$ comme voulu.
\end{proof}

\begin{remark}
\label{remark-negative-dimension}
Notons que, dans le contexte du lemme précédent,
on n'a pas nécessairement $\dim \mathcal{T} \geq \dim \mathcal{Z}$; cela ne
contredit pas l'inégalité de l'énoncé du lemme, car
les fibres du morphisme $f$ sont encore des champs algébriques, et
peuvent donc avoir une dimension négative. On l'illustre en prenant
pour $k$ un corps et en appliquant le lemme au morphisme
$[\Spec k/\mathbf{G}_m] \to \Spec k$.

\medskip\noindent
Si le morphisme $f$ de l'énoncé du lemme est supposé
quasi-DM (au sens de
Morphismes de champs, Définition
\ref{stacks-morphisms-definition-separated}; les morphismes
représentables par des espaces algébriques sont par exemple quasi-DM),
alors les fibres du morphisme aux points du but
sont des champs algébriques quasi-DM, et ont donc une dimension
positive ou nulle. Dans ce cas, le lemme entraîne
que l'on a bien $\dim \mathcal{T} \geq \dim \mathcal{Z}$. En fait, nous obtenons
le résultat plus général suivant.
\end{remark}

\begin{lemma}
\label{lemma-dims-of-images-two}
Soit $f: \mathcal{T} \to \mathcal{X}$ un morphisme localement de type fini
de champs algébriques jacobsoniens, pseudo-caténaires et localement
noethériens,
qui est quasi-DM,
dont la source est irréductible et le but quasi-séparé,
et soit $\mathcal{Z} \hookrightarrow \mathcal{X}$ l'image
schématique de $\mathcal{T}$.
Alors $\dim \mathcal{Z} \leq \dim \mathcal{T}$,
et, de plus, exactement l'une des deux conditions suivantes est satisfaite :
\begin{enumerate}
\item pour tout point de type fini $t \in |T|,$
nous avons
$\dim_t(\mathcal{T}_{f(t)}) > 0,$ auquel
cas $\dim \mathcal{Z} < \dim \mathcal{T}$; ou
\item $\mathcal{T}$ et $\mathcal{Z}$
ont la même dimension.
\end{enumerate}
\end{lemma}

\begin{proof}
Comme nous l'avons observé dans la remarque précédente,
la dimension d'un champ quasi-DM est toujours positive ou nulle,
d'où nous concluons que $\dim_t \mathcal{T}_{f(t)} \geq 0$
pour tout $t \in |\mathcal{T}|$, l'égalité
$$
\dim_t \mathcal{T}_{f(t)} = \dim_t \mathcal{T} - \dim_{f(t)} \mathcal{Z}
$$
étant satisfaite
sur un sous-ensemble ouvert dense de points $t\in |\mathcal{T}|$.
\end{proof}





\section{La dimension de l'anneau local}
\label{section-dimension-local-ring}

\noindent
Un champ algébrique ne possède pas véritablement d'anneaux locaux au sens usuel,
mais nous pouvons définir la dimension de l'anneau local comme suit.

\begin{lemma}
\label{lemma-dimension-local-ring-pre}
Soit $\mathcal{X}$ un champ algébrique localement noethérien.
Soit $U \to \mathcal{X}$ un morphisme lisse et soit $u \in U$.
Alors
$$
\dim(\mathcal{O}_{U, \overline{u}}) -
\dim(\mathcal{O}_{R_u, e(\overline{u})}) =
2\dim(\mathcal{O}_{U, \overline{u}}) -
\dim(\mathcal{O}_{R, e(\overline{u})})
$$
Ici, $R = U \times_\mathcal{X} U$, de projections $s, t : R \to U$ et de
diagonale $e : U \to R$, et $R_u$ est la fibre de $s : R \to U$ en $u$.
\end{lemma}

\begin{proof}
Cela est vrai parce que
$s : \mathcal{O}_{U, \overline{u}} \to \mathcal{O}_{R, e(\overline{u})}$
est un homomorphisme local plat d'anneaux locaux noethériens, et donc
$$
\dim(\mathcal{O}_{R, e(\overline{u})}) =
\dim(\mathcal{O}_{U, \overline{u}}) +
\dim(\mathcal{O}_{R_u, e(\overline{u})})
$$
d'après
Algèbre, Lemme \ref{algebra-lemma-dimension-base-fibre-equals-total}.
\end{proof}

\begin{lemma}
\label{lemma-dimension-local-ring}
Soit $\mathcal{X}$ un champ algébrique localement noethérien.
Soit $x \in |\mathcal{X}|$ un point de type fini 
Morphismes de champs, Définition
\ref{stacks-morphisms-definition-finite-type-point}).
Soit $d \in \mathbf{Z}$.
Les conditions suivantes sont équivalentes :
\begin{enumerate}
\item il existe un schéma $U$, un morphisme lisse $U \to \mathcal{X}$,
et un point de type fini $u \in U$ s'envoyant sur $x$, tels que
$2\dim(\mathcal{O}_{U, \overline{u}}) -
\dim(\mathcal{O}_{R, e(\overline{u})}) = d$, et
\item pour tout schéma $U$, tout morphisme lisse $U \to \mathcal{X}$,
et tout point de type fini $u \in U$ s'envoyant sur $x$, nous avons
$2\dim(\mathcal{O}_{U, \overline{u}}) -
\dim(\mathcal{O}_{R, e(\overline{u})}) = d$.
\end{enumerate}
Ici, $R = U \times_\mathcal{X} U$, de projections $s, t : R \to U$ et de
diagonale $e : U \to R$, et $R_u$ est la fibre de $s : R \to U$ en $u$.
\end{lemma}

\begin{proof}
Supposons donnés deux voisinages lisses $(U, u)$ et $(U', u')$
de $x$, où $u$ et $u'$ sont des points de type fini. Quitte à rétrécir $U$
et $U'$, nous pouvons supposer que $u$ et $u'$ sont des points fermés
(par définition des points de type fini). Choisissons alors
un morphisme étale surjectif $W \to U \times_\mathcal{X} U'$.
Soit $W_u$ la fibre de $W \to U$ en $u$, et
soit $W_{u'}$ la fibre de $W \to U'$ en $u'$.
Puisque $u$ et $u'$ s'envoient sur le même point de $|\mathcal{X}|$,
nous voyons que $W_u \cap W_{u'}$ est non vide.
Nous pouvons donc choisir un point fermé $w \in W$ s'envoyant à la fois
sur $u$ et sur $u'$. Cela nous ramène à la discussion du
paragraphe suivant.

\medskip\noindent
Supposons que $(U', u') \to (U, u)$ soit un morphisme lisse de voisinages
lisses de $x$, avec $u$ et $u'$ des points fermés.
Le but est de démontrer que l'invariant défini pour $(U, u)$ est égal
à celui défini pour $(U', u')$.
Pour le voir, observons que $\mathcal{O}_{U, u} \to \mathcal{O}_{U', u'}$
est un homomorphisme local plat d'anneaux locaux noethériens, et donc
$$
\dim(\mathcal{O}_{U', \overline{u}'}) =
\dim(\mathcal{O}_{U, \overline{u}}) +
\dim(\mathcal{O}_{U'_u, \overline{u}'})
$$
d'après Algèbre, Lemme \ref{algebra-lemma-dimension-base-fibre-equals-total}.
(Nous omettons le détail de toutes les étapes reliant les propriétés des anneaux
locaux et de leurs hensélisés stricts; voir
Compléments d'algèbre, section \ref{more-algebra-section-permanence-henselization}).
D'autre part, nous avons
$$
R' = U' \times_{U, t} R \times_{s, U} U'
$$
Nous voyons donc que
$$
\dim(\mathcal{O}_{R', e(\overline{u}')}) =
\dim(\mathcal{O}_{R, e(\overline{u})}) +
\dim(\mathcal{O}_{U'_u \times_u U'_u, (\overline{u}', \overline{u}')})
$$
Pour démontrer le lemme, il suffit de montrer que
$$
\dim(\mathcal{O}_{U'_u \times_u U'_u, (\overline{u}', \overline{u}')}) =
2\dim(\mathcal{O}_{U'_u, \overline{u}'})
$$
Observons que ceci n'est pas toujours vrai (par exemple, si $U'_u$ est une courbe
et $u'$ le point générique de cette courbe). Nous savons toutefois
que $u'$ est un point fermé de l'espace algébrique $U'_u$, localement
de type fini sur $u$. Dans ce cas, l'égalité est satisfaite car, premièrement,
$\dim_{(u', u')}(U'_u \times_u U'_u) = 2\dim_{u'}(U'_u)$ d'après
Variétés, Lemme \ref{varieties-lemma-dimension-product-locally-algebraic},
et, deuxièmement, parce que la dimension coïncide avec celle des anneaux locaux
aux points fermés des schémas localement algébriques; voir
Variétés, Lemme \ref{varieties-lemma-dimension-locally-algebraic}.
Nous omettons la transposition de ces résultats pour les schémas dans le
langage des espaces algébriques.
\end{proof}

\begin{definition}
\label{definition-dimension-local-ring}
Soit $\mathcal{X}$ un champ algébrique localement noethérien.
Soit $x \in |\mathcal{X}|$ un point de type fini.
La {\it dimension de l'anneau local de $\mathcal{X}$ en $x$}
est $d \in \mathbf{Z}$ si les conditions équivalentes du
Lemme \ref{lemma-dimension-local-ring} sont satisfaites.
\end{definition}

\noindent
Bien entendu, ceci est motivé par le Lemme \ref{lemma-dimension-local-ring-pre}
et Propriétés des champs, Définition
\ref{stacks-properties-definition-dimension-at-point}.
Terminons cette section en établissant une formule permettant de
calculer $\dim_x(\mathcal{X})$ à partir des propriétés de l'anneau versel
de $\mathcal{X}$ en $x$.

\begin{lemma}
\label{lemma-dimension-formula}
Supposons que $\mathcal{X}$ soit un champ algébrique, localement de type fini
sur un schéma localement noethérien $S$. Soit $x_0 : \Spec(k) \to \mathcal{X}$
un morphisme, où $k$ est un corps de type fini sur $S$. Représentons
$\mathcal{F}_{\mathcal{X}, k, x_0}$ comme dans la Remarque \ref{remark-groupoid-defo}
par un cogroupoïde $(A, B, s, t, c)$ de $S$-algèbres locales noethériennes complètes
de corps résiduel $k$. Alors
$$
\text{la dimension de l'anneau local de }\mathcal{X}\text{ en }x_0 =
2\dim A - \dim B
$$
\end{lemma}

\begin{proof}
Soit $s \in S$ l'image de $x_0$. Si $\mathcal{O}_{S, s}$
est un anneau G (condition presque toujours satisfaite en pratique),
nous pouvons démontrer le lemme comme suit.
D'après le Lemme \ref{lemma-Artin-approximation-by-smooth-morphism},
nous pouvons trouver un morphisme lisse $U \to \mathcal{X}$, de source un schéma,
contenant un point $u_0 \in U$ de corps résiduel $k$, tel que le
morphisme induit $\Spec(k) \to U \to \mathcal{X}$ coïncide avec $x_0$
et que $A = \mathcal{O}_{U, u_0}^\wedge$.
Écrivons $R = U \times_\mathcal{X} U$. Nous pouvons alors identifier
$\mathcal{O}_{R, e(u_0)}^\wedge$ à $B$.
L'égalité résulte donc des définitions.

\medskip\noindent
Dans la suite de cette preuve, nous expliquons comment démontrer le lemme en
général, mais nous invitons vivement le lecteur à passer ce qui suit.

\medskip\noindent
Montrons d'abord que le membre de droite est indépendant du choix
de $(A, B, s, t, c)$. Supposons en effet que $(A', B', s', t', c')$
soit un second choix. Puisque $A$ et $A'$ sont des anneaux versels de $\mathcal{X}$
en $x_0$, nous pouvons choisir, quitte à échanger $A$ et $A'$,
une application formellement lisse $A \to A'$ compatible avec les objets formels
versels donnés $\xi$ et $\xi'$ sur $A$ et $A'$.
Rappelons que $\widehat{\mathcal{C}}_\Lambda$ possède
des coproduits, donnés par le produit tensoriel complété
sur $\Lambda$; voir Théorie formelle des déformations, Lemme
\ref{formal-defos-lemma-CLambdahat-coproducts}.
Alors $B$ proreprésente le foncteur des isomorphismes entre les
deux images directes de $\xi$ dans $A \widehat{\otimes}_\Lambda A$.
Il en va de même pour $B'$. Nous en concluons que
$$
B' =
B \otimes_{(A \widehat{\otimes}_\Lambda A)}
(A' \widehat{\otimes}_\Lambda A')
$$
On voit immédiatement que
$$
A \widehat{\otimes}_\Lambda A \longrightarrow
A \widehat{\otimes}_\Lambda A' \longrightarrow
A' \widehat{\otimes}_\Lambda A'
$$
est formellement lisse, de dimension relative égale à $2$ fois la
dimension relative de l'application formellement lisse $A \to A'$.
(Cela résulte de principes généraux, mais aussi du fait que,
dans ce cas particulier, $A'$ est un anneau de séries formelles sur $A$
en $r$ variables.) Ainsi, $B \to B'$ est formellement lisse, de
dimension relative $2(\dim(A') - \dim(A))$, comme voulu.

\medskip\noindent
Ensuite, soit $l/k$ une extension finie. Soit
$x_{l, 0} : \Spec(l) \to \mathcal{X}$
le point induit. Nous affirmons que le membre de droite de la formule
est le même pour $x_0$ que pour $x_{l, 0}$.
On le montre en choisissant $A \to A'$ comme dans le
Lemme \ref{lemma-versal-ring-field-extension}
et en raisonnant exactement comme au paragraphe précédent.
Nous omettons les détails.

\medskip\noindent
Enfin, en raisonnant comme dans la preuve du Lemme \ref{lemma-versal-ring-flat},
nous pouvons utiliser les compatibilités des deux paragraphes précédents
pour nous ramener au cas (étudié au premier paragraphe)
où $A$ est l'anneau local complet de $U$ en $u_0$, pour un certain
schéma lisse sur $\mathcal{X}$ et un point de type fini $u_0$.
Détails omis.
\end{proof}






\begin{multicols}{2}[\section{Autres chapitres}]
\noindent
Préliminaires
\begin{enumerate}
\item \hyperref[introduction-section-phantom]{Introduction}
\item \hyperref[conventions-section-phantom]{Conventions}
\item \hyperref[sets-section-phantom]{Théorie des ensembles}
\item \hyperref[categories-section-phantom]{Catégories}
\item \hyperref[topology-section-phantom]{Topologie}
\item \hyperref[sheaves-section-phantom]{Faisceaux sur les espaces}
\item \hyperref[sites-section-phantom]{Sites et faisceaux}
\item \hyperref[stacks-section-phantom]{Champs}
\item \hyperref[fields-section-phantom]{Corps}
\item \hyperref[algebra-section-phantom]{Algèbre commutative}
\item \hyperref[brauer-section-phantom]{Groupes de Brauer}
\item \hyperref[homology-section-phantom]{Algèbre homologique}
\item \hyperref[derived-section-phantom]{Catégories dérivées}
\item \hyperref[simplicial-section-phantom]{Méthodes simpliciales}
\item \hyperref[more-algebra-section-phantom]{Compléments d'algèbre}
\item \hyperref[smoothing-section-phantom]{Lissage des morphismes d'anneaux}
\item \hyperref[modules-section-phantom]{Faisceaux de modules}
\item \hyperref[sites-modules-section-phantom]{Modules sur les sites}
\item \hyperref[injectives-section-phantom]{Objets injectifs}
\item \hyperref[cohomology-section-phantom]{Cohomologie des faisceaux}
\item \hyperref[sites-cohomology-section-phantom]{Cohomologie sur les sites}
\item \hyperref[dga-section-phantom]{Algèbre différentielle graduée}
\item \hyperref[dpa-section-phantom]{Algèbre à puissances divisées}
\item \hyperref[sdga-section-phantom]{Faisceaux différentiels gradués}
\item \hyperref[hypercovering-section-phantom]{Hyperrecouvrements}
\end{enumerate}
Schémas
\begin{enumerate}
\setcounter{enumi}{25}
\item \hyperref[schemes-section-phantom]{Schémas}
\item \hyperref[constructions-section-phantom]{Constructions de schémas}
\item \hyperref[properties-section-phantom]{Propriétés des schémas}
\item \hyperref[morphisms-section-phantom]{Morphismes de schémas}
\item \hyperref[coherent-section-phantom]{Cohomologie des schémas}
\item \hyperref[divisors-section-phantom]{Diviseurs}
\item \hyperref[limits-section-phantom]{Limites de schémas}
\item \hyperref[varieties-section-phantom]{Variétés}
\item \hyperref[topologies-section-phantom]{Topologies sur les schémas}
\item \hyperref[descent-section-phantom]{Descente}
\item \hyperref[perfect-section-phantom]{Catégories dérivées de schémas}
\item \hyperref[more-morphisms-section-phantom]{Compléments sur les morphismes}
\item \hyperref[flat-section-phantom]{Compléments sur la platitude}
\item \hyperref[groupoids-section-phantom]{Schémas en groupoïdes}
\item \hyperref[more-groupoids-section-phantom]{Compléments sur les schémas en groupoïdes}
\item \hyperref[etale-section-phantom]{Morphismes étales de schémas}
\end{enumerate}
Thèmes de théorie des schémas
\begin{enumerate}
\setcounter{enumi}{41}
\item \hyperref[chow-section-phantom]{Homologie de Chow}
\item \hyperref[intersection-section-phantom]{Théorie de l'intersection}
\item \hyperref[pic-section-phantom]{Schémas de Picard des courbes}
\item \hyperref[weil-section-phantom]{Théories cohomologiques de Weil}
\item \hyperref[adequate-section-phantom]{Modules adéquats}
\item \hyperref[dualizing-section-phantom]{Complexes dualisants}
\item \hyperref[duality-section-phantom]{Dualité pour les schémas}
\item \hyperref[discriminant-section-phantom]{Discriminants et différentes}
\item \hyperref[derham-section-phantom]{Cohomologie de de Rham}
\item \hyperref[local-cohomology-section-phantom]{Cohomologie locale}
\item \hyperref[algebraization-section-phantom]{Géométrie algébrique et formelle}
\item \hyperref[curves-section-phantom]{Courbes algébriques}
\item \hyperref[resolve-section-phantom]{Résolution des surfaces}
\item \hyperref[models-section-phantom]{Réduction semi-stable}
\item \hyperref[functors-section-phantom]{Foncteurs et morphismes}
\item \hyperref[equiv-section-phantom]{Catégories dérivées de variétés}
\item \hyperref[pione-section-phantom]{Groupes fondamentaux de schémas}
\item \hyperref[etale-cohomology-section-phantom]{Cohomologie étale}
\item \hyperref[crystalline-section-phantom]{Cohomologie cristalline}
\item \hyperref[proetale-section-phantom]{Cohomologie pro-étale}
\item \hyperref[relative-cycles-section-phantom]{Cycles relatifs}
\item \hyperref[more-etale-section-phantom]{Compléments de cohomologie étale}
\item \hyperref[trace-section-phantom]{La formule des traces}
\end{enumerate}
Espaces algébriques
\begin{enumerate}
\setcounter{enumi}{64}
\item \hyperref[spaces-section-phantom]{Espaces algébriques}
\item \hyperref[spaces-properties-section-phantom]{Propriétés des espaces algébriques}
\item \hyperref[spaces-morphisms-section-phantom]{Morphismes d'espaces algébriques}
\item \hyperref[decent-spaces-section-phantom]{Espaces algébriques décents}
\item \hyperref[spaces-cohomology-section-phantom]{Cohomologie des espaces algébriques}
\item \hyperref[spaces-limits-section-phantom]{Limites d'espaces algébriques}
\item \hyperref[spaces-divisors-section-phantom]{Diviseurs sur les espaces algébriques}
\item \hyperref[spaces-over-fields-section-phantom]{Espaces algébriques sur des corps}
\item \hyperref[spaces-topologies-section-phantom]{Topologies sur les espaces algébriques}
\item \hyperref[spaces-descent-section-phantom]{Descente et espaces algébriques}
\item \hyperref[spaces-perfect-section-phantom]{Catégories dérivées d'espaces}
\item \hyperref[spaces-more-morphisms-section-phantom]{Compléments sur les morphismes d'espaces}
\item \hyperref[spaces-flat-section-phantom]{Platitude sur les espaces algébriques}
\item \hyperref[spaces-groupoids-section-phantom]{Groupoïdes dans les espaces algébriques}
\item \hyperref[spaces-more-groupoids-section-phantom]{Compléments sur les groupoïdes dans les espaces}
\item \hyperref[bootstrap-section-phantom]{Amorçage}
\item \hyperref[spaces-pushouts-section-phantom]{Sommes amalgamées d'espaces algébriques}
\end{enumerate}
Thèmes de géométrie
\begin{enumerate}
\setcounter{enumi}{81}
\item \hyperref[spaces-chow-section-phantom]{Groupes de Chow des espaces}
\item \hyperref[groupoids-quotients-section-phantom]{Quotients de groupoïdes}
\item \hyperref[spaces-more-cohomology-section-phantom]{Compléments sur la cohomologie des espaces}
\item \hyperref[spaces-simplicial-section-phantom]{Espaces simpliciaux}
\item \hyperref[spaces-duality-section-phantom]{Dualité pour les espaces}
\item \hyperref[formal-spaces-section-phantom]{Espaces algébriques formels}
\item \hyperref[restricted-section-phantom]{Algébrisation des espaces formels}
\item \hyperref[spaces-resolve-section-phantom]{Résolution des surfaces revisitée}
\end{enumerate}
Théorie des déformations
\begin{enumerate}
\setcounter{enumi}{89}
\item \hyperref[formal-defos-section-phantom]{Théorie formelle des déformations}
\item \hyperref[defos-section-phantom]{Théorie des déformations}
\item \hyperref[cotangent-section-phantom]{Le complexe cotangent}
\item \hyperref[examples-defos-section-phantom]{Problèmes de déformation}
\end{enumerate}
Champs algébriques
\begin{enumerate}
\setcounter{enumi}{93}
\item \hyperref[algebraic-section-phantom]{Champs algébriques}
\item \hyperref[examples-stacks-section-phantom]{Exemples de champs}
\item \hyperref[stacks-sheaves-section-phantom]{Faisceaux sur les champs algébriques}
\item \hyperref[criteria-section-phantom]{Critères de représentabilité}
\item \hyperref[artin-section-phantom]{Axiomes d'Artin}
\item \hyperref[quot-section-phantom]{Espaces de Quot et de Hilbert}
\item \hyperref[stacks-properties-section-phantom]{Propriétés des champs algébriques}
\item \hyperref[stacks-morphisms-section-phantom]{Morphismes de champs algébriques}
\item \hyperref[stacks-limits-section-phantom]{Limites de champs algébriques}
\item \hyperref[stacks-cohomology-section-phantom]{Cohomologie des champs algébriques}
\item \hyperref[stacks-perfect-section-phantom]{Catégories dérivées de champs}
\item \hyperref[stacks-introduction-section-phantom]{Introduction aux champs algébriques}
\item \hyperref[stacks-more-morphisms-section-phantom]{Compléments sur les morphismes de champs}
\item \hyperref[stacks-geometry-section-phantom]{La géométrie des champs}
\end{enumerate}
Thèmes de théorie des espaces de modules
\begin{enumerate}
\setcounter{enumi}{107}
\item \hyperref[moduli-section-phantom]{Champs de modules}
\item \hyperref[moduli-curves-section-phantom]{Espaces de modules de courbes}
\end{enumerate}
Divers
\begin{enumerate}
\setcounter{enumi}{109}
\item \hyperref[examples-section-phantom]{Exemples}
\item \hyperref[exercises-section-phantom]{Exercices}
\item \hyperref[guide-section-phantom]{Guide bibliographique}
\item \hyperref[desirables-section-phantom]{Desiderata}
\item \hyperref[coding-section-phantom]{Style de programmation}
\item \hyperref[obsolete-section-phantom]{Obsolète}
\item \hyperref[fdl-section-phantom]{Licence de documentation libre GNU}
\item \hyperref[index-section-phantom]{Index généré automatiquement}
\end{enumerate}
\end{multicols}


\bibliography{my}
\bibliographystyle{amsalpha}

\end{document}
